% ============================================================
% 论文带读 PDF 统一主题模板（XeLaTeX + ctexart）
% 使用：% ============================================================
% 论文带读 PDF 统一主题模板（XeLaTeX + ctexart）
% 使用：% ============================================================
% 论文带读 PDF 统一主题模板（XeLaTeX + ctexart）
% 使用：% ============================================================
% 论文带读 PDF 统一主题模板（XeLaTeX + ctexart）
% 使用：\input{preamble}
% 注意：不要用 \accent 作命令名（与 LaTeX 内置冲突），用 \highlight。
% ============================================================
\documentclass[11pt,a4paper]{ctexart}

\usepackage[a4paper,margin=19mm,headheight=23pt,headsep=7mm,footskip=10mm]{geometry}
\usepackage{fontspec}
\setmainfont{Times New Roman}
\setsansfont{Microsoft YaHei}
\setmonofont{Consolas}[Scale=0.84]
\usepackage{amsmath,amssymb}
\usepackage{graphicx}
\usepackage{booktabs,tabularx,array,multirow}
\usepackage{caption,subcaption}
\usepackage{longtable}
\usepackage{makecell}
\usepackage{enumitem}
\usepackage{float}
\usepackage{listings}
\usepackage{tcolorbox}
\tcbuselibrary{skins,breakable}
\usepackage{tikz}
\usetikzlibrary{arrows.meta,positioning,shapes.geometric,fit,backgrounds}
\usepackage{fancyhdr}
\usepackage{lastpage}
\usepackage{hyperref}
\usepackage{xcolor}
\usepackage{titlesec}
\usepackage{setspace}
\usepackage{ragged2e}
\usepackage{pifont}

\definecolor{DeepTeal}{HTML}{0E5A67}
\definecolor{MidTeal}{HTML}{3E8C88}
\definecolor{WarmOrange}{HTML}{D17A35}
\definecolor{Ink}{HTML}{1D2B30}
\definecolor{PaleBlue}{HTML}{F3F7F8}
\definecolor{PaleTeal}{HTML}{E7F2F1}
\definecolor{CodeBg}{HTML}{F5F7F7}
\definecolor{CodeRule}{HTML}{C6D5D6}

% 建议在正文 \input{preamble} 之后用 \hypersetup 覆盖标题
\hypersetup{colorlinks=true,linkcolor=DeepTeal,urlcolor=DeepTeal,citecolor=DeepTeal,
  pdftitle={论文带读},pdfauthor={论文带读}}
\captionsetup{font=small,labelfont={bf,color=DeepTeal},textfont={color=Ink},skip=5pt}
\setlength{\parindent}{2em}
\setlength{\parskip}{0.28em}
\linespread{1.12}

\titleformat{\section}{\Large\bfseries\color{DeepTeal}}{}{0pt}{}
\titlespacing*{\section}{0pt}{0pt}{6pt}
\titleformat{\subsection}{\large\bfseries\color{DeepTeal}}{}{0pt}{}
\titlespacing*{\subsection}{0pt}{5pt}{3pt}

\pagestyle{fancy}
\fancyhf{}
\lhead{\small\color{DeepTeal}\textbf{文献阅读汇报}\quad <论文主题>}
\rhead{\small\color{gray}<作者>, <年>}
\cfoot{\small\color{gray}\thepage\,/\,\pageref{LastPage}}
\renewcommand{\headrulewidth}{0.5pt}
\renewcommand{\headrule}{\hbox to\headwidth{\color{DeepTeal}\leaders\hrule height \headrulewidth\hfill}}

% ---------- 版式组件 ----------
\newcommand{\key}[1]{\textcolor{DeepTeal}{\textbf{#1}}}
\newcommand{\highlight}[1]{\textcolor{WarmOrange}{\textbf{#1}}}
\newcommand{\source}[1]{\par\vspace{1pt}\raggedright\scriptsize\color{gray}\textit{来源：#1}\par}
\newcommand{\sectiontag}[1]{\noindent\colorbox{PaleTeal}{\parbox{\dimexpr\linewidth-2\fboxsep\relax}{\small\color{DeepTeal}\textbf{#1}}}\par\vspace{5pt}}

% 本节路线（浅青绿框）
\newcommand{\roadmap}[1]{\begin{tcolorbox}[colback=PaleTeal,colframe=MidTeal,boxrule=.5pt,arc=1.5mm,left=2.2mm,right=2.2mm,top=1.4mm,bottom=1.4mm]\small\textbf{本节路线：}\quad #1\end{tcolorbox}}
% 大白话（浅蓝框）
\newcommand{\plain}[1]{\begin{tcolorbox}[colback=PaleBlue,colframe=DeepTeal,boxrule=.5pt,arc=1.5mm,left=2.2mm,right=2.2mm,top=1.5mm,bottom=1.5mm]\small\textbf{大白话}\quad #1\end{tcolorbox}}
% 专业定义 + 本文中的位置（白底青绿框）
\newcommand{\definition}[2]{\begin{tcolorbox}[colback=white,colframe=MidTeal,boxrule=.45pt,arc=1.5mm,left=2.2mm,right=2.2mm,top=1.4mm,bottom=1.4mm]\small\textbf{专业定义：}#1\par\vspace{1mm}\textbf{本文中的位置：}#2\end{tcolorbox}}
% 不要这样理解（暖橙框）
\newcommand{\warning}[1]{\begin{tcolorbox}[colback=WarmOrange!10,colframe=WarmOrange,boxrule=.5pt,arc=1.5mm,left=2.2mm,right=2.2mm,top=1.4mm,bottom=1.4mm]\small\textbf{不要这样理解：}\quad #1\end{tcolorbox}}
% 阅读要点（青绿框）
\newcommand{\takeaway}[1]{\begin{tcolorbox}[colback=PaleTeal,colframe=MidTeal,boxrule=.55pt,arc=1.5mm,left=2.2mm,right=2.2mm,top=1.5mm,bottom=1.5mm]\small\textbf{阅读要点}\quad #1\end{tcolorbox}}
% 注意（白底暖橙框）
\newcommand{\noteBox}[1]{\begin{tcolorbox}[colback=white,colframe=WarmOrange,boxrule=.55pt,arc=1.5mm,left=2.2mm,right=2.2mm,top=1.3mm,bottom=1.3mm]\small #1\end{tcolorbox}}
% 论文图（紧跟图注）
\newcommand{\paperfig}[3]{\begin{figure}[H]\centering\includegraphics[width=#2\linewidth]{#1}\caption{#3}\end{figure}}
% 代码文件标注
\newcommand{\codefile}[2]{\noindent\scriptsize\color{gray}\texttt{#1}\quad\textit{#2}\par\vspace{2pt}}

% ---------- 代码样式 ----------
\lstdefinestyle{pythonstyle}{language=Python,basicstyle=\ttfamily\scriptsize\color{Ink},
  keywordstyle=\color{DeepTeal}\bfseries,commentstyle=\color{gray}\itshape,
  stringstyle=\color{WarmOrange},backgroundcolor=\color{CodeBg},frame=single,
  rulecolor=\color{CodeRule},framerule=.35pt,breaklines=true,showstringspaces=false,
  columns=fullflexible,aboveskip=3pt,belowskip=3pt,xleftmargin=1pt,xrightmargin=1pt}

% ---------- 表格列类型 ----------
\newcolumntype{Y}{>{\RaggedRight\arraybackslash}X}
\newcolumntype{C}{>{\Centering\arraybackslash}X}

% ---------- 安全表格（避免 p{}+Y 混用溢出；见 SKILL.md 陷阱 6） ----------
% 用法：
%   \begin{center}\small
%   \begin{tabular}{@{}p{0.26\linewidth}@{\hspace{4mm}}p{0.62\linewidth}@{}}
%   ...\end{tabular}
%   \end{center}

% 注意：不要用 \accent 作命令名（与 LaTeX 内置冲突），用 \highlight。
% ============================================================
\documentclass[11pt,a4paper]{ctexart}

\usepackage[a4paper,margin=19mm,headheight=23pt,headsep=7mm,footskip=10mm]{geometry}
\usepackage{fontspec}
\setmainfont{Times New Roman}
\setsansfont{Microsoft YaHei}
\setmonofont{Consolas}[Scale=0.84]
\usepackage{amsmath,amssymb}
\usepackage{graphicx}
\usepackage{booktabs,tabularx,array,multirow}
\usepackage{caption,subcaption}
\usepackage{longtable}
\usepackage{makecell}
\usepackage{enumitem}
\usepackage{float}
\usepackage{listings}
\usepackage{tcolorbox}
\tcbuselibrary{skins,breakable}
\usepackage{tikz}
\usetikzlibrary{arrows.meta,positioning,shapes.geometric,fit,backgrounds}
\usepackage{fancyhdr}
\usepackage{lastpage}
\usepackage{hyperref}
\usepackage{xcolor}
\usepackage{titlesec}
\usepackage{setspace}
\usepackage{ragged2e}
\usepackage{pifont}

\definecolor{DeepTeal}{HTML}{0E5A67}
\definecolor{MidTeal}{HTML}{3E8C88}
\definecolor{WarmOrange}{HTML}{D17A35}
\definecolor{Ink}{HTML}{1D2B30}
\definecolor{PaleBlue}{HTML}{F3F7F8}
\definecolor{PaleTeal}{HTML}{E7F2F1}
\definecolor{CodeBg}{HTML}{F5F7F7}
\definecolor{CodeRule}{HTML}{C6D5D6}

% 建议在正文 % ============================================================
% 论文带读 PDF 统一主题模板（XeLaTeX + ctexart）
% 使用：\input{preamble}
% 注意：不要用 \accent 作命令名（与 LaTeX 内置冲突），用 \highlight。
% ============================================================
\documentclass[11pt,a4paper]{ctexart}

\usepackage[a4paper,margin=19mm,headheight=23pt,headsep=7mm,footskip=10mm]{geometry}
\usepackage{fontspec}
\setmainfont{Times New Roman}
\setsansfont{Microsoft YaHei}
\setmonofont{Consolas}[Scale=0.84]
\usepackage{amsmath,amssymb}
\usepackage{graphicx}
\usepackage{booktabs,tabularx,array,multirow}
\usepackage{caption,subcaption}
\usepackage{longtable}
\usepackage{makecell}
\usepackage{enumitem}
\usepackage{float}
\usepackage{listings}
\usepackage{tcolorbox}
\tcbuselibrary{skins,breakable}
\usepackage{tikz}
\usetikzlibrary{arrows.meta,positioning,shapes.geometric,fit,backgrounds}
\usepackage{fancyhdr}
\usepackage{lastpage}
\usepackage{hyperref}
\usepackage{xcolor}
\usepackage{titlesec}
\usepackage{setspace}
\usepackage{ragged2e}
\usepackage{pifont}

\definecolor{DeepTeal}{HTML}{0E5A67}
\definecolor{MidTeal}{HTML}{3E8C88}
\definecolor{WarmOrange}{HTML}{D17A35}
\definecolor{Ink}{HTML}{1D2B30}
\definecolor{PaleBlue}{HTML}{F3F7F8}
\definecolor{PaleTeal}{HTML}{E7F2F1}
\definecolor{CodeBg}{HTML}{F5F7F7}
\definecolor{CodeRule}{HTML}{C6D5D6}

% 建议在正文 \input{preamble} 之后用 \hypersetup 覆盖标题
\hypersetup{colorlinks=true,linkcolor=DeepTeal,urlcolor=DeepTeal,citecolor=DeepTeal,
  pdftitle={论文带读},pdfauthor={论文带读}}
\captionsetup{font=small,labelfont={bf,color=DeepTeal},textfont={color=Ink},skip=5pt}
\setlength{\parindent}{2em}
\setlength{\parskip}{0.28em}
\linespread{1.12}

\titleformat{\section}{\Large\bfseries\color{DeepTeal}}{}{0pt}{}
\titlespacing*{\section}{0pt}{0pt}{6pt}
\titleformat{\subsection}{\large\bfseries\color{DeepTeal}}{}{0pt}{}
\titlespacing*{\subsection}{0pt}{5pt}{3pt}

\pagestyle{fancy}
\fancyhf{}
\lhead{\small\color{DeepTeal}\textbf{文献阅读汇报}\quad <论文主题>}
\rhead{\small\color{gray}<作者>, <年>}
\cfoot{\small\color{gray}\thepage\,/\,\pageref{LastPage}}
\renewcommand{\headrulewidth}{0.5pt}
\renewcommand{\headrule}{\hbox to\headwidth{\color{DeepTeal}\leaders\hrule height \headrulewidth\hfill}}

% ---------- 版式组件 ----------
\newcommand{\key}[1]{\textcolor{DeepTeal}{\textbf{#1}}}
\newcommand{\highlight}[1]{\textcolor{WarmOrange}{\textbf{#1}}}
\newcommand{\source}[1]{\par\vspace{1pt}\raggedright\scriptsize\color{gray}\textit{来源：#1}\par}
\newcommand{\sectiontag}[1]{\noindent\colorbox{PaleTeal}{\parbox{\dimexpr\linewidth-2\fboxsep\relax}{\small\color{DeepTeal}\textbf{#1}}}\par\vspace{5pt}}

% 本节路线（浅青绿框）
\newcommand{\roadmap}[1]{\begin{tcolorbox}[colback=PaleTeal,colframe=MidTeal,boxrule=.5pt,arc=1.5mm,left=2.2mm,right=2.2mm,top=1.4mm,bottom=1.4mm]\small\textbf{本节路线：}\quad #1\end{tcolorbox}}
% 大白话（浅蓝框）
\newcommand{\plain}[1]{\begin{tcolorbox}[colback=PaleBlue,colframe=DeepTeal,boxrule=.5pt,arc=1.5mm,left=2.2mm,right=2.2mm,top=1.5mm,bottom=1.5mm]\small\textbf{大白话}\quad #1\end{tcolorbox}}
% 专业定义 + 本文中的位置（白底青绿框）
\newcommand{\definition}[2]{\begin{tcolorbox}[colback=white,colframe=MidTeal,boxrule=.45pt,arc=1.5mm,left=2.2mm,right=2.2mm,top=1.4mm,bottom=1.4mm]\small\textbf{专业定义：}#1\par\vspace{1mm}\textbf{本文中的位置：}#2\end{tcolorbox}}
% 不要这样理解（暖橙框）
\newcommand{\warning}[1]{\begin{tcolorbox}[colback=WarmOrange!10,colframe=WarmOrange,boxrule=.5pt,arc=1.5mm,left=2.2mm,right=2.2mm,top=1.4mm,bottom=1.4mm]\small\textbf{不要这样理解：}\quad #1\end{tcolorbox}}
% 阅读要点（青绿框）
\newcommand{\takeaway}[1]{\begin{tcolorbox}[colback=PaleTeal,colframe=MidTeal,boxrule=.55pt,arc=1.5mm,left=2.2mm,right=2.2mm,top=1.5mm,bottom=1.5mm]\small\textbf{阅读要点}\quad #1\end{tcolorbox}}
% 注意（白底暖橙框）
\newcommand{\noteBox}[1]{\begin{tcolorbox}[colback=white,colframe=WarmOrange,boxrule=.55pt,arc=1.5mm,left=2.2mm,right=2.2mm,top=1.3mm,bottom=1.3mm]\small #1\end{tcolorbox}}
% 论文图（紧跟图注）
\newcommand{\paperfig}[3]{\begin{figure}[H]\centering\includegraphics[width=#2\linewidth]{#1}\caption{#3}\end{figure}}
% 代码文件标注
\newcommand{\codefile}[2]{\noindent\scriptsize\color{gray}\texttt{#1}\quad\textit{#2}\par\vspace{2pt}}

% ---------- 代码样式 ----------
\lstdefinestyle{pythonstyle}{language=Python,basicstyle=\ttfamily\scriptsize\color{Ink},
  keywordstyle=\color{DeepTeal}\bfseries,commentstyle=\color{gray}\itshape,
  stringstyle=\color{WarmOrange},backgroundcolor=\color{CodeBg},frame=single,
  rulecolor=\color{CodeRule},framerule=.35pt,breaklines=true,showstringspaces=false,
  columns=fullflexible,aboveskip=3pt,belowskip=3pt,xleftmargin=1pt,xrightmargin=1pt}

% ---------- 表格列类型 ----------
\newcolumntype{Y}{>{\RaggedRight\arraybackslash}X}
\newcolumntype{C}{>{\Centering\arraybackslash}X}

% ---------- 安全表格（避免 p{}+Y 混用溢出；见 SKILL.md 陷阱 6） ----------
% 用法：
%   \begin{center}\small
%   \begin{tabular}{@{}p{0.26\linewidth}@{\hspace{4mm}}p{0.62\linewidth}@{}}
%   ...\end{tabular}
%   \end{center}
 之后用 \hypersetup 覆盖标题
\hypersetup{colorlinks=true,linkcolor=DeepTeal,urlcolor=DeepTeal,citecolor=DeepTeal,
  pdftitle={论文带读},pdfauthor={论文带读}}
\captionsetup{font=small,labelfont={bf,color=DeepTeal},textfont={color=Ink},skip=5pt}
\setlength{\parindent}{2em}
\setlength{\parskip}{0.28em}
\linespread{1.12}

\titleformat{\section}{\Large\bfseries\color{DeepTeal}}{}{0pt}{}
\titlespacing*{\section}{0pt}{0pt}{6pt}
\titleformat{\subsection}{\large\bfseries\color{DeepTeal}}{}{0pt}{}
\titlespacing*{\subsection}{0pt}{5pt}{3pt}

\pagestyle{fancy}
\fancyhf{}
\lhead{\small\color{DeepTeal}\textbf{文献阅读汇报}\quad <论文主题>}
\rhead{\small\color{gray}<作者>, <年>}
\cfoot{\small\color{gray}\thepage\,/\,\pageref{LastPage}}
\renewcommand{\headrulewidth}{0.5pt}
\renewcommand{\headrule}{\hbox to\headwidth{\color{DeepTeal}\leaders\hrule height \headrulewidth\hfill}}

% ---------- 版式组件 ----------
\newcommand{\key}[1]{\textcolor{DeepTeal}{\textbf{#1}}}
\newcommand{\highlight}[1]{\textcolor{WarmOrange}{\textbf{#1}}}
\newcommand{\source}[1]{\par\vspace{1pt}\raggedright\scriptsize\color{gray}\textit{来源：#1}\par}
\newcommand{\sectiontag}[1]{\noindent\colorbox{PaleTeal}{\parbox{\dimexpr\linewidth-2\fboxsep\relax}{\small\color{DeepTeal}\textbf{#1}}}\par\vspace{5pt}}

% 本节路线（浅青绿框）
\newcommand{\roadmap}[1]{\begin{tcolorbox}[colback=PaleTeal,colframe=MidTeal,boxrule=.5pt,arc=1.5mm,left=2.2mm,right=2.2mm,top=1.4mm,bottom=1.4mm]\small\textbf{本节路线：}\quad #1\end{tcolorbox}}
% 大白话（浅蓝框）
\newcommand{\plain}[1]{\begin{tcolorbox}[colback=PaleBlue,colframe=DeepTeal,boxrule=.5pt,arc=1.5mm,left=2.2mm,right=2.2mm,top=1.5mm,bottom=1.5mm]\small\textbf{大白话}\quad #1\end{tcolorbox}}
% 专业定义 + 本文中的位置（白底青绿框）
\newcommand{\definition}[2]{\begin{tcolorbox}[colback=white,colframe=MidTeal,boxrule=.45pt,arc=1.5mm,left=2.2mm,right=2.2mm,top=1.4mm,bottom=1.4mm]\small\textbf{专业定义：}#1\par\vspace{1mm}\textbf{本文中的位置：}#2\end{tcolorbox}}
% 不要这样理解（暖橙框）
\newcommand{\warning}[1]{\begin{tcolorbox}[colback=WarmOrange!10,colframe=WarmOrange,boxrule=.5pt,arc=1.5mm,left=2.2mm,right=2.2mm,top=1.4mm,bottom=1.4mm]\small\textbf{不要这样理解：}\quad #1\end{tcolorbox}}
% 阅读要点（青绿框）
\newcommand{\takeaway}[1]{\begin{tcolorbox}[colback=PaleTeal,colframe=MidTeal,boxrule=.55pt,arc=1.5mm,left=2.2mm,right=2.2mm,top=1.5mm,bottom=1.5mm]\small\textbf{阅读要点}\quad #1\end{tcolorbox}}
% 注意（白底暖橙框）
\newcommand{\noteBox}[1]{\begin{tcolorbox}[colback=white,colframe=WarmOrange,boxrule=.55pt,arc=1.5mm,left=2.2mm,right=2.2mm,top=1.3mm,bottom=1.3mm]\small #1\end{tcolorbox}}
% 论文图（紧跟图注）
\newcommand{\paperfig}[3]{\begin{figure}[H]\centering\includegraphics[width=#2\linewidth]{#1}\caption{#3}\end{figure}}
% 代码文件标注
\newcommand{\codefile}[2]{\noindent\scriptsize\color{gray}\texttt{#1}\quad\textit{#2}\par\vspace{2pt}}

% ---------- 代码样式 ----------
\lstdefinestyle{pythonstyle}{language=Python,basicstyle=\ttfamily\scriptsize\color{Ink},
  keywordstyle=\color{DeepTeal}\bfseries,commentstyle=\color{gray}\itshape,
  stringstyle=\color{WarmOrange},backgroundcolor=\color{CodeBg},frame=single,
  rulecolor=\color{CodeRule},framerule=.35pt,breaklines=true,showstringspaces=false,
  columns=fullflexible,aboveskip=3pt,belowskip=3pt,xleftmargin=1pt,xrightmargin=1pt}

% ---------- 表格列类型 ----------
\newcolumntype{Y}{>{\RaggedRight\arraybackslash}X}
\newcolumntype{C}{>{\Centering\arraybackslash}X}

% ---------- 安全表格（避免 p{}+Y 混用溢出；见 SKILL.md 陷阱 6） ----------
% 用法：
%   \begin{center}\small
%   \begin{tabular}{@{}p{0.26\linewidth}@{\hspace{4mm}}p{0.62\linewidth}@{}}
%   ...\end{tabular}
%   \end{center}

% 注意：不要用 \accent 作命令名（与 LaTeX 内置冲突），用 \highlight。
% ============================================================
\documentclass[11pt,a4paper]{ctexart}

\usepackage[a4paper,margin=19mm,headheight=23pt,headsep=7mm,footskip=10mm]{geometry}
\usepackage{fontspec}
\setmainfont{Times New Roman}
\setsansfont{Microsoft YaHei}
\setmonofont{Consolas}[Scale=0.84]
\usepackage{amsmath,amssymb}
\usepackage{graphicx}
\usepackage{booktabs,tabularx,array,multirow}
\usepackage{caption,subcaption}
\usepackage{longtable}
\usepackage{makecell}
\usepackage{enumitem}
\usepackage{float}
\usepackage{listings}
\usepackage{tcolorbox}
\tcbuselibrary{skins,breakable}
\usepackage{tikz}
\usetikzlibrary{arrows.meta,positioning,shapes.geometric,fit,backgrounds}
\usepackage{fancyhdr}
\usepackage{lastpage}
\usepackage{hyperref}
\usepackage{xcolor}
\usepackage{titlesec}
\usepackage{setspace}
\usepackage{ragged2e}
\usepackage{pifont}

\definecolor{DeepTeal}{HTML}{0E5A67}
\definecolor{MidTeal}{HTML}{3E8C88}
\definecolor{WarmOrange}{HTML}{D17A35}
\definecolor{Ink}{HTML}{1D2B30}
\definecolor{PaleBlue}{HTML}{F3F7F8}
\definecolor{PaleTeal}{HTML}{E7F2F1}
\definecolor{CodeBg}{HTML}{F5F7F7}
\definecolor{CodeRule}{HTML}{C6D5D6}

% 建议在正文 % ============================================================
% 论文带读 PDF 统一主题模板（XeLaTeX + ctexart）
% 使用：% ============================================================
% 论文带读 PDF 统一主题模板（XeLaTeX + ctexart）
% 使用：\input{preamble}
% 注意：不要用 \accent 作命令名（与 LaTeX 内置冲突），用 \highlight。
% ============================================================
\documentclass[11pt,a4paper]{ctexart}

\usepackage[a4paper,margin=19mm,headheight=23pt,headsep=7mm,footskip=10mm]{geometry}
\usepackage{fontspec}
\setmainfont{Times New Roman}
\setsansfont{Microsoft YaHei}
\setmonofont{Consolas}[Scale=0.84]
\usepackage{amsmath,amssymb}
\usepackage{graphicx}
\usepackage{booktabs,tabularx,array,multirow}
\usepackage{caption,subcaption}
\usepackage{longtable}
\usepackage{makecell}
\usepackage{enumitem}
\usepackage{float}
\usepackage{listings}
\usepackage{tcolorbox}
\tcbuselibrary{skins,breakable}
\usepackage{tikz}
\usetikzlibrary{arrows.meta,positioning,shapes.geometric,fit,backgrounds}
\usepackage{fancyhdr}
\usepackage{lastpage}
\usepackage{hyperref}
\usepackage{xcolor}
\usepackage{titlesec}
\usepackage{setspace}
\usepackage{ragged2e}
\usepackage{pifont}

\definecolor{DeepTeal}{HTML}{0E5A67}
\definecolor{MidTeal}{HTML}{3E8C88}
\definecolor{WarmOrange}{HTML}{D17A35}
\definecolor{Ink}{HTML}{1D2B30}
\definecolor{PaleBlue}{HTML}{F3F7F8}
\definecolor{PaleTeal}{HTML}{E7F2F1}
\definecolor{CodeBg}{HTML}{F5F7F7}
\definecolor{CodeRule}{HTML}{C6D5D6}

% 建议在正文 \input{preamble} 之后用 \hypersetup 覆盖标题
\hypersetup{colorlinks=true,linkcolor=DeepTeal,urlcolor=DeepTeal,citecolor=DeepTeal,
  pdftitle={论文带读},pdfauthor={论文带读}}
\captionsetup{font=small,labelfont={bf,color=DeepTeal},textfont={color=Ink},skip=5pt}
\setlength{\parindent}{2em}
\setlength{\parskip}{0.28em}
\linespread{1.12}

\titleformat{\section}{\Large\bfseries\color{DeepTeal}}{}{0pt}{}
\titlespacing*{\section}{0pt}{0pt}{6pt}
\titleformat{\subsection}{\large\bfseries\color{DeepTeal}}{}{0pt}{}
\titlespacing*{\subsection}{0pt}{5pt}{3pt}

\pagestyle{fancy}
\fancyhf{}
\lhead{\small\color{DeepTeal}\textbf{文献阅读汇报}\quad <论文主题>}
\rhead{\small\color{gray}<作者>, <年>}
\cfoot{\small\color{gray}\thepage\,/\,\pageref{LastPage}}
\renewcommand{\headrulewidth}{0.5pt}
\renewcommand{\headrule}{\hbox to\headwidth{\color{DeepTeal}\leaders\hrule height \headrulewidth\hfill}}

% ---------- 版式组件 ----------
\newcommand{\key}[1]{\textcolor{DeepTeal}{\textbf{#1}}}
\newcommand{\highlight}[1]{\textcolor{WarmOrange}{\textbf{#1}}}
\newcommand{\source}[1]{\par\vspace{1pt}\raggedright\scriptsize\color{gray}\textit{来源：#1}\par}
\newcommand{\sectiontag}[1]{\noindent\colorbox{PaleTeal}{\parbox{\dimexpr\linewidth-2\fboxsep\relax}{\small\color{DeepTeal}\textbf{#1}}}\par\vspace{5pt}}

% 本节路线（浅青绿框）
\newcommand{\roadmap}[1]{\begin{tcolorbox}[colback=PaleTeal,colframe=MidTeal,boxrule=.5pt,arc=1.5mm,left=2.2mm,right=2.2mm,top=1.4mm,bottom=1.4mm]\small\textbf{本节路线：}\quad #1\end{tcolorbox}}
% 大白话（浅蓝框）
\newcommand{\plain}[1]{\begin{tcolorbox}[colback=PaleBlue,colframe=DeepTeal,boxrule=.5pt,arc=1.5mm,left=2.2mm,right=2.2mm,top=1.5mm,bottom=1.5mm]\small\textbf{大白话}\quad #1\end{tcolorbox}}
% 专业定义 + 本文中的位置（白底青绿框）
\newcommand{\definition}[2]{\begin{tcolorbox}[colback=white,colframe=MidTeal,boxrule=.45pt,arc=1.5mm,left=2.2mm,right=2.2mm,top=1.4mm,bottom=1.4mm]\small\textbf{专业定义：}#1\par\vspace{1mm}\textbf{本文中的位置：}#2\end{tcolorbox}}
% 不要这样理解（暖橙框）
\newcommand{\warning}[1]{\begin{tcolorbox}[colback=WarmOrange!10,colframe=WarmOrange,boxrule=.5pt,arc=1.5mm,left=2.2mm,right=2.2mm,top=1.4mm,bottom=1.4mm]\small\textbf{不要这样理解：}\quad #1\end{tcolorbox}}
% 阅读要点（青绿框）
\newcommand{\takeaway}[1]{\begin{tcolorbox}[colback=PaleTeal,colframe=MidTeal,boxrule=.55pt,arc=1.5mm,left=2.2mm,right=2.2mm,top=1.5mm,bottom=1.5mm]\small\textbf{阅读要点}\quad #1\end{tcolorbox}}
% 注意（白底暖橙框）
\newcommand{\noteBox}[1]{\begin{tcolorbox}[colback=white,colframe=WarmOrange,boxrule=.55pt,arc=1.5mm,left=2.2mm,right=2.2mm,top=1.3mm,bottom=1.3mm]\small #1\end{tcolorbox}}
% 论文图（紧跟图注）
\newcommand{\paperfig}[3]{\begin{figure}[H]\centering\includegraphics[width=#2\linewidth]{#1}\caption{#3}\end{figure}}
% 代码文件标注
\newcommand{\codefile}[2]{\noindent\scriptsize\color{gray}\texttt{#1}\quad\textit{#2}\par\vspace{2pt}}

% ---------- 代码样式 ----------
\lstdefinestyle{pythonstyle}{language=Python,basicstyle=\ttfamily\scriptsize\color{Ink},
  keywordstyle=\color{DeepTeal}\bfseries,commentstyle=\color{gray}\itshape,
  stringstyle=\color{WarmOrange},backgroundcolor=\color{CodeBg},frame=single,
  rulecolor=\color{CodeRule},framerule=.35pt,breaklines=true,showstringspaces=false,
  columns=fullflexible,aboveskip=3pt,belowskip=3pt,xleftmargin=1pt,xrightmargin=1pt}

% ---------- 表格列类型 ----------
\newcolumntype{Y}{>{\RaggedRight\arraybackslash}X}
\newcolumntype{C}{>{\Centering\arraybackslash}X}

% ---------- 安全表格（避免 p{}+Y 混用溢出；见 SKILL.md 陷阱 6） ----------
% 用法：
%   \begin{center}\small
%   \begin{tabular}{@{}p{0.26\linewidth}@{\hspace{4mm}}p{0.62\linewidth}@{}}
%   ...\end{tabular}
%   \end{center}

% 注意：不要用 \accent 作命令名（与 LaTeX 内置冲突），用 \highlight。
% ============================================================
\documentclass[11pt,a4paper]{ctexart}

\usepackage[a4paper,margin=19mm,headheight=23pt,headsep=7mm,footskip=10mm]{geometry}
\usepackage{fontspec}
\setmainfont{Times New Roman}
\setsansfont{Microsoft YaHei}
\setmonofont{Consolas}[Scale=0.84]
\usepackage{amsmath,amssymb}
\usepackage{graphicx}
\usepackage{booktabs,tabularx,array,multirow}
\usepackage{caption,subcaption}
\usepackage{longtable}
\usepackage{makecell}
\usepackage{enumitem}
\usepackage{float}
\usepackage{listings}
\usepackage{tcolorbox}
\tcbuselibrary{skins,breakable}
\usepackage{tikz}
\usetikzlibrary{arrows.meta,positioning,shapes.geometric,fit,backgrounds}
\usepackage{fancyhdr}
\usepackage{lastpage}
\usepackage{hyperref}
\usepackage{xcolor}
\usepackage{titlesec}
\usepackage{setspace}
\usepackage{ragged2e}
\usepackage{pifont}

\definecolor{DeepTeal}{HTML}{0E5A67}
\definecolor{MidTeal}{HTML}{3E8C88}
\definecolor{WarmOrange}{HTML}{D17A35}
\definecolor{Ink}{HTML}{1D2B30}
\definecolor{PaleBlue}{HTML}{F3F7F8}
\definecolor{PaleTeal}{HTML}{E7F2F1}
\definecolor{CodeBg}{HTML}{F5F7F7}
\definecolor{CodeRule}{HTML}{C6D5D6}

% 建议在正文 % ============================================================
% 论文带读 PDF 统一主题模板（XeLaTeX + ctexart）
% 使用：\input{preamble}
% 注意：不要用 \accent 作命令名（与 LaTeX 内置冲突），用 \highlight。
% ============================================================
\documentclass[11pt,a4paper]{ctexart}

\usepackage[a4paper,margin=19mm,headheight=23pt,headsep=7mm,footskip=10mm]{geometry}
\usepackage{fontspec}
\setmainfont{Times New Roman}
\setsansfont{Microsoft YaHei}
\setmonofont{Consolas}[Scale=0.84]
\usepackage{amsmath,amssymb}
\usepackage{graphicx}
\usepackage{booktabs,tabularx,array,multirow}
\usepackage{caption,subcaption}
\usepackage{longtable}
\usepackage{makecell}
\usepackage{enumitem}
\usepackage{float}
\usepackage{listings}
\usepackage{tcolorbox}
\tcbuselibrary{skins,breakable}
\usepackage{tikz}
\usetikzlibrary{arrows.meta,positioning,shapes.geometric,fit,backgrounds}
\usepackage{fancyhdr}
\usepackage{lastpage}
\usepackage{hyperref}
\usepackage{xcolor}
\usepackage{titlesec}
\usepackage{setspace}
\usepackage{ragged2e}
\usepackage{pifont}

\definecolor{DeepTeal}{HTML}{0E5A67}
\definecolor{MidTeal}{HTML}{3E8C88}
\definecolor{WarmOrange}{HTML}{D17A35}
\definecolor{Ink}{HTML}{1D2B30}
\definecolor{PaleBlue}{HTML}{F3F7F8}
\definecolor{PaleTeal}{HTML}{E7F2F1}
\definecolor{CodeBg}{HTML}{F5F7F7}
\definecolor{CodeRule}{HTML}{C6D5D6}

% 建议在正文 \input{preamble} 之后用 \hypersetup 覆盖标题
\hypersetup{colorlinks=true,linkcolor=DeepTeal,urlcolor=DeepTeal,citecolor=DeepTeal,
  pdftitle={论文带读},pdfauthor={论文带读}}
\captionsetup{font=small,labelfont={bf,color=DeepTeal},textfont={color=Ink},skip=5pt}
\setlength{\parindent}{2em}
\setlength{\parskip}{0.28em}
\linespread{1.12}

\titleformat{\section}{\Large\bfseries\color{DeepTeal}}{}{0pt}{}
\titlespacing*{\section}{0pt}{0pt}{6pt}
\titleformat{\subsection}{\large\bfseries\color{DeepTeal}}{}{0pt}{}
\titlespacing*{\subsection}{0pt}{5pt}{3pt}

\pagestyle{fancy}
\fancyhf{}
\lhead{\small\color{DeepTeal}\textbf{文献阅读汇报}\quad <论文主题>}
\rhead{\small\color{gray}<作者>, <年>}
\cfoot{\small\color{gray}\thepage\,/\,\pageref{LastPage}}
\renewcommand{\headrulewidth}{0.5pt}
\renewcommand{\headrule}{\hbox to\headwidth{\color{DeepTeal}\leaders\hrule height \headrulewidth\hfill}}

% ---------- 版式组件 ----------
\newcommand{\key}[1]{\textcolor{DeepTeal}{\textbf{#1}}}
\newcommand{\highlight}[1]{\textcolor{WarmOrange}{\textbf{#1}}}
\newcommand{\source}[1]{\par\vspace{1pt}\raggedright\scriptsize\color{gray}\textit{来源：#1}\par}
\newcommand{\sectiontag}[1]{\noindent\colorbox{PaleTeal}{\parbox{\dimexpr\linewidth-2\fboxsep\relax}{\small\color{DeepTeal}\textbf{#1}}}\par\vspace{5pt}}

% 本节路线（浅青绿框）
\newcommand{\roadmap}[1]{\begin{tcolorbox}[colback=PaleTeal,colframe=MidTeal,boxrule=.5pt,arc=1.5mm,left=2.2mm,right=2.2mm,top=1.4mm,bottom=1.4mm]\small\textbf{本节路线：}\quad #1\end{tcolorbox}}
% 大白话（浅蓝框）
\newcommand{\plain}[1]{\begin{tcolorbox}[colback=PaleBlue,colframe=DeepTeal,boxrule=.5pt,arc=1.5mm,left=2.2mm,right=2.2mm,top=1.5mm,bottom=1.5mm]\small\textbf{大白话}\quad #1\end{tcolorbox}}
% 专业定义 + 本文中的位置（白底青绿框）
\newcommand{\definition}[2]{\begin{tcolorbox}[colback=white,colframe=MidTeal,boxrule=.45pt,arc=1.5mm,left=2.2mm,right=2.2mm,top=1.4mm,bottom=1.4mm]\small\textbf{专业定义：}#1\par\vspace{1mm}\textbf{本文中的位置：}#2\end{tcolorbox}}
% 不要这样理解（暖橙框）
\newcommand{\warning}[1]{\begin{tcolorbox}[colback=WarmOrange!10,colframe=WarmOrange,boxrule=.5pt,arc=1.5mm,left=2.2mm,right=2.2mm,top=1.4mm,bottom=1.4mm]\small\textbf{不要这样理解：}\quad #1\end{tcolorbox}}
% 阅读要点（青绿框）
\newcommand{\takeaway}[1]{\begin{tcolorbox}[colback=PaleTeal,colframe=MidTeal,boxrule=.55pt,arc=1.5mm,left=2.2mm,right=2.2mm,top=1.5mm,bottom=1.5mm]\small\textbf{阅读要点}\quad #1\end{tcolorbox}}
% 注意（白底暖橙框）
\newcommand{\noteBox}[1]{\begin{tcolorbox}[colback=white,colframe=WarmOrange,boxrule=.55pt,arc=1.5mm,left=2.2mm,right=2.2mm,top=1.3mm,bottom=1.3mm]\small #1\end{tcolorbox}}
% 论文图（紧跟图注）
\newcommand{\paperfig}[3]{\begin{figure}[H]\centering\includegraphics[width=#2\linewidth]{#1}\caption{#3}\end{figure}}
% 代码文件标注
\newcommand{\codefile}[2]{\noindent\scriptsize\color{gray}\texttt{#1}\quad\textit{#2}\par\vspace{2pt}}

% ---------- 代码样式 ----------
\lstdefinestyle{pythonstyle}{language=Python,basicstyle=\ttfamily\scriptsize\color{Ink},
  keywordstyle=\color{DeepTeal}\bfseries,commentstyle=\color{gray}\itshape,
  stringstyle=\color{WarmOrange},backgroundcolor=\color{CodeBg},frame=single,
  rulecolor=\color{CodeRule},framerule=.35pt,breaklines=true,showstringspaces=false,
  columns=fullflexible,aboveskip=3pt,belowskip=3pt,xleftmargin=1pt,xrightmargin=1pt}

% ---------- 表格列类型 ----------
\newcolumntype{Y}{>{\RaggedRight\arraybackslash}X}
\newcolumntype{C}{>{\Centering\arraybackslash}X}

% ---------- 安全表格（避免 p{}+Y 混用溢出；见 SKILL.md 陷阱 6） ----------
% 用法：
%   \begin{center}\small
%   \begin{tabular}{@{}p{0.26\linewidth}@{\hspace{4mm}}p{0.62\linewidth}@{}}
%   ...\end{tabular}
%   \end{center}
 之后用 \hypersetup 覆盖标题
\hypersetup{colorlinks=true,linkcolor=DeepTeal,urlcolor=DeepTeal,citecolor=DeepTeal,
  pdftitle={论文带读},pdfauthor={论文带读}}
\captionsetup{font=small,labelfont={bf,color=DeepTeal},textfont={color=Ink},skip=5pt}
\setlength{\parindent}{2em}
\setlength{\parskip}{0.28em}
\linespread{1.12}

\titleformat{\section}{\Large\bfseries\color{DeepTeal}}{}{0pt}{}
\titlespacing*{\section}{0pt}{0pt}{6pt}
\titleformat{\subsection}{\large\bfseries\color{DeepTeal}}{}{0pt}{}
\titlespacing*{\subsection}{0pt}{5pt}{3pt}

\pagestyle{fancy}
\fancyhf{}
\lhead{\small\color{DeepTeal}\textbf{文献阅读汇报}\quad <论文主题>}
\rhead{\small\color{gray}<作者>, <年>}
\cfoot{\small\color{gray}\thepage\,/\,\pageref{LastPage}}
\renewcommand{\headrulewidth}{0.5pt}
\renewcommand{\headrule}{\hbox to\headwidth{\color{DeepTeal}\leaders\hrule height \headrulewidth\hfill}}

% ---------- 版式组件 ----------
\newcommand{\key}[1]{\textcolor{DeepTeal}{\textbf{#1}}}
\newcommand{\highlight}[1]{\textcolor{WarmOrange}{\textbf{#1}}}
\newcommand{\source}[1]{\par\vspace{1pt}\raggedright\scriptsize\color{gray}\textit{来源：#1}\par}
\newcommand{\sectiontag}[1]{\noindent\colorbox{PaleTeal}{\parbox{\dimexpr\linewidth-2\fboxsep\relax}{\small\color{DeepTeal}\textbf{#1}}}\par\vspace{5pt}}

% 本节路线（浅青绿框）
\newcommand{\roadmap}[1]{\begin{tcolorbox}[colback=PaleTeal,colframe=MidTeal,boxrule=.5pt,arc=1.5mm,left=2.2mm,right=2.2mm,top=1.4mm,bottom=1.4mm]\small\textbf{本节路线：}\quad #1\end{tcolorbox}}
% 大白话（浅蓝框）
\newcommand{\plain}[1]{\begin{tcolorbox}[colback=PaleBlue,colframe=DeepTeal,boxrule=.5pt,arc=1.5mm,left=2.2mm,right=2.2mm,top=1.5mm,bottom=1.5mm]\small\textbf{大白话}\quad #1\end{tcolorbox}}
% 专业定义 + 本文中的位置（白底青绿框）
\newcommand{\definition}[2]{\begin{tcolorbox}[colback=white,colframe=MidTeal,boxrule=.45pt,arc=1.5mm,left=2.2mm,right=2.2mm,top=1.4mm,bottom=1.4mm]\small\textbf{专业定义：}#1\par\vspace{1mm}\textbf{本文中的位置：}#2\end{tcolorbox}}
% 不要这样理解（暖橙框）
\newcommand{\warning}[1]{\begin{tcolorbox}[colback=WarmOrange!10,colframe=WarmOrange,boxrule=.5pt,arc=1.5mm,left=2.2mm,right=2.2mm,top=1.4mm,bottom=1.4mm]\small\textbf{不要这样理解：}\quad #1\end{tcolorbox}}
% 阅读要点（青绿框）
\newcommand{\takeaway}[1]{\begin{tcolorbox}[colback=PaleTeal,colframe=MidTeal,boxrule=.55pt,arc=1.5mm,left=2.2mm,right=2.2mm,top=1.5mm,bottom=1.5mm]\small\textbf{阅读要点}\quad #1\end{tcolorbox}}
% 注意（白底暖橙框）
\newcommand{\noteBox}[1]{\begin{tcolorbox}[colback=white,colframe=WarmOrange,boxrule=.55pt,arc=1.5mm,left=2.2mm,right=2.2mm,top=1.3mm,bottom=1.3mm]\small #1\end{tcolorbox}}
% 论文图（紧跟图注）
\newcommand{\paperfig}[3]{\begin{figure}[H]\centering\includegraphics[width=#2\linewidth]{#1}\caption{#3}\end{figure}}
% 代码文件标注
\newcommand{\codefile}[2]{\noindent\scriptsize\color{gray}\texttt{#1}\quad\textit{#2}\par\vspace{2pt}}

% ---------- 代码样式 ----------
\lstdefinestyle{pythonstyle}{language=Python,basicstyle=\ttfamily\scriptsize\color{Ink},
  keywordstyle=\color{DeepTeal}\bfseries,commentstyle=\color{gray}\itshape,
  stringstyle=\color{WarmOrange},backgroundcolor=\color{CodeBg},frame=single,
  rulecolor=\color{CodeRule},framerule=.35pt,breaklines=true,showstringspaces=false,
  columns=fullflexible,aboveskip=3pt,belowskip=3pt,xleftmargin=1pt,xrightmargin=1pt}

% ---------- 表格列类型 ----------
\newcolumntype{Y}{>{\RaggedRight\arraybackslash}X}
\newcolumntype{C}{>{\Centering\arraybackslash}X}

% ---------- 安全表格（避免 p{}+Y 混用溢出；见 SKILL.md 陷阱 6） ----------
% 用法：
%   \begin{center}\small
%   \begin{tabular}{@{}p{0.26\linewidth}@{\hspace{4mm}}p{0.62\linewidth}@{}}
%   ...\end{tabular}
%   \end{center}
 之后用 \hypersetup 覆盖标题
\hypersetup{colorlinks=true,linkcolor=DeepTeal,urlcolor=DeepTeal,citecolor=DeepTeal,
  pdftitle={论文带读},pdfauthor={论文带读}}
\captionsetup{font=small,labelfont={bf,color=DeepTeal},textfont={color=Ink},skip=5pt}
\setlength{\parindent}{2em}
\setlength{\parskip}{0.28em}
\linespread{1.12}

\titleformat{\section}{\Large\bfseries\color{DeepTeal}}{}{0pt}{}
\titlespacing*{\section}{0pt}{0pt}{6pt}
\titleformat{\subsection}{\large\bfseries\color{DeepTeal}}{}{0pt}{}
\titlespacing*{\subsection}{0pt}{5pt}{3pt}

\pagestyle{fancy}
\fancyhf{}
\lhead{\small\color{DeepTeal}\textbf{文献阅读汇报}\quad <论文主题>}
\rhead{\small\color{gray}<作者>, <年>}
\cfoot{\small\color{gray}\thepage\,/\,\pageref{LastPage}}
\renewcommand{\headrulewidth}{0.5pt}
\renewcommand{\headrule}{\hbox to\headwidth{\color{DeepTeal}\leaders\hrule height \headrulewidth\hfill}}

% ---------- 版式组件 ----------
\newcommand{\key}[1]{\textcolor{DeepTeal}{\textbf{#1}}}
\newcommand{\highlight}[1]{\textcolor{WarmOrange}{\textbf{#1}}}
\newcommand{\source}[1]{\par\vspace{1pt}\raggedright\scriptsize\color{gray}\textit{来源：#1}\par}
\newcommand{\sectiontag}[1]{\noindent\colorbox{PaleTeal}{\parbox{\dimexpr\linewidth-2\fboxsep\relax}{\small\color{DeepTeal}\textbf{#1}}}\par\vspace{5pt}}

% 本节路线（浅青绿框）
\newcommand{\roadmap}[1]{\begin{tcolorbox}[colback=PaleTeal,colframe=MidTeal,boxrule=.5pt,arc=1.5mm,left=2.2mm,right=2.2mm,top=1.4mm,bottom=1.4mm]\small\textbf{本节路线：}\quad #1\end{tcolorbox}}
% 大白话（浅蓝框）
\newcommand{\plain}[1]{\begin{tcolorbox}[colback=PaleBlue,colframe=DeepTeal,boxrule=.5pt,arc=1.5mm,left=2.2mm,right=2.2mm,top=1.5mm,bottom=1.5mm]\small\textbf{大白话}\quad #1\end{tcolorbox}}
% 专业定义 + 本文中的位置（白底青绿框）
\newcommand{\definition}[2]{\begin{tcolorbox}[colback=white,colframe=MidTeal,boxrule=.45pt,arc=1.5mm,left=2.2mm,right=2.2mm,top=1.4mm,bottom=1.4mm]\small\textbf{专业定义：}#1\par\vspace{1mm}\textbf{本文中的位置：}#2\end{tcolorbox}}
% 不要这样理解（暖橙框）
\newcommand{\warning}[1]{\begin{tcolorbox}[colback=WarmOrange!10,colframe=WarmOrange,boxrule=.5pt,arc=1.5mm,left=2.2mm,right=2.2mm,top=1.4mm,bottom=1.4mm]\small\textbf{不要这样理解：}\quad #1\end{tcolorbox}}
% 阅读要点（青绿框）
\newcommand{\takeaway}[1]{\begin{tcolorbox}[colback=PaleTeal,colframe=MidTeal,boxrule=.55pt,arc=1.5mm,left=2.2mm,right=2.2mm,top=1.5mm,bottom=1.5mm]\small\textbf{阅读要点}\quad #1\end{tcolorbox}}
% 注意（白底暖橙框）
\newcommand{\noteBox}[1]{\begin{tcolorbox}[colback=white,colframe=WarmOrange,boxrule=.55pt,arc=1.5mm,left=2.2mm,right=2.2mm,top=1.3mm,bottom=1.3mm]\small #1\end{tcolorbox}}
% 论文图（紧跟图注）
\newcommand{\paperfig}[3]{\begin{figure}[H]\centering\includegraphics[width=#2\linewidth]{#1}\caption{#3}\end{figure}}
% 代码文件标注
\newcommand{\codefile}[2]{\noindent\scriptsize\color{gray}\texttt{#1}\quad\textit{#2}\par\vspace{2pt}}

% ---------- 代码样式 ----------
\lstdefinestyle{pythonstyle}{language=Python,basicstyle=\ttfamily\scriptsize\color{Ink},
  keywordstyle=\color{DeepTeal}\bfseries,commentstyle=\color{gray}\itshape,
  stringstyle=\color{WarmOrange},backgroundcolor=\color{CodeBg},frame=single,
  rulecolor=\color{CodeRule},framerule=.35pt,breaklines=true,showstringspaces=false,
  columns=fullflexible,aboveskip=3pt,belowskip=3pt,xleftmargin=1pt,xrightmargin=1pt}

% ---------- 表格列类型 ----------
\newcolumntype{Y}{>{\RaggedRight\arraybackslash}X}
\newcolumntype{C}{>{\Centering\arraybackslash}X}

% ---------- 安全表格（避免 p{}+Y 混用溢出；见 SKILL.md 陷阱 6） ----------
% 用法：
%   \begin{center}\small
%   \begin{tabular}{@{}p{0.26\linewidth}@{\hspace{4mm}}p{0.62\linewidth}@{}}
%   ...\end{tabular}
%   \end{center}

% 注意：不要用 \accent 作命令名（与 LaTeX 内置冲突），用 \highlight。
% ============================================================
\documentclass[11pt,a4paper]{ctexart}

\usepackage[a4paper,margin=19mm,headheight=23pt,headsep=7mm,footskip=10mm]{geometry}
\usepackage{fontspec}
\setmainfont{Times New Roman}
\setsansfont{Microsoft YaHei}
\setmonofont{Consolas}[Scale=0.84]
\usepackage{amsmath,amssymb}
\usepackage{graphicx}
\usepackage{booktabs,tabularx,array,multirow}
\usepackage{caption,subcaption}
\usepackage{longtable}
\usepackage{makecell}
\usepackage{enumitem}
\usepackage{float}
\usepackage{listings}
\usepackage{tcolorbox}
\tcbuselibrary{skins,breakable}
\usepackage{tikz}
\usetikzlibrary{arrows.meta,positioning,shapes.geometric,fit,backgrounds}
\usepackage{fancyhdr}
\usepackage{lastpage}
\usepackage{hyperref}
\usepackage{xcolor}
\usepackage{titlesec}
\usepackage{setspace}
\usepackage{ragged2e}
\usepackage{pifont}

\definecolor{DeepTeal}{HTML}{0E5A67}
\definecolor{MidTeal}{HTML}{3E8C88}
\definecolor{WarmOrange}{HTML}{D17A35}
\definecolor{Ink}{HTML}{1D2B30}
\definecolor{PaleBlue}{HTML}{F3F7F8}
\definecolor{PaleTeal}{HTML}{E7F2F1}
\definecolor{CodeBg}{HTML}{F5F7F7}
\definecolor{CodeRule}{HTML}{C6D5D6}

% 建议在正文 % ============================================================
% 论文带读 PDF 统一主题模板（XeLaTeX + ctexart）
% 使用：% ============================================================
% 论文带读 PDF 统一主题模板（XeLaTeX + ctexart）
% 使用：% ============================================================
% 论文带读 PDF 统一主题模板（XeLaTeX + ctexart）
% 使用：\input{preamble}
% 注意：不要用 \accent 作命令名（与 LaTeX 内置冲突），用 \highlight。
% ============================================================
\documentclass[11pt,a4paper]{ctexart}

\usepackage[a4paper,margin=19mm,headheight=23pt,headsep=7mm,footskip=10mm]{geometry}
\usepackage{fontspec}
\setmainfont{Times New Roman}
\setsansfont{Microsoft YaHei}
\setmonofont{Consolas}[Scale=0.84]
\usepackage{amsmath,amssymb}
\usepackage{graphicx}
\usepackage{booktabs,tabularx,array,multirow}
\usepackage{caption,subcaption}
\usepackage{longtable}
\usepackage{makecell}
\usepackage{enumitem}
\usepackage{float}
\usepackage{listings}
\usepackage{tcolorbox}
\tcbuselibrary{skins,breakable}
\usepackage{tikz}
\usetikzlibrary{arrows.meta,positioning,shapes.geometric,fit,backgrounds}
\usepackage{fancyhdr}
\usepackage{lastpage}
\usepackage{hyperref}
\usepackage{xcolor}
\usepackage{titlesec}
\usepackage{setspace}
\usepackage{ragged2e}
\usepackage{pifont}

\definecolor{DeepTeal}{HTML}{0E5A67}
\definecolor{MidTeal}{HTML}{3E8C88}
\definecolor{WarmOrange}{HTML}{D17A35}
\definecolor{Ink}{HTML}{1D2B30}
\definecolor{PaleBlue}{HTML}{F3F7F8}
\definecolor{PaleTeal}{HTML}{E7F2F1}
\definecolor{CodeBg}{HTML}{F5F7F7}
\definecolor{CodeRule}{HTML}{C6D5D6}

% 建议在正文 \input{preamble} 之后用 \hypersetup 覆盖标题
\hypersetup{colorlinks=true,linkcolor=DeepTeal,urlcolor=DeepTeal,citecolor=DeepTeal,
  pdftitle={论文带读},pdfauthor={论文带读}}
\captionsetup{font=small,labelfont={bf,color=DeepTeal},textfont={color=Ink},skip=5pt}
\setlength{\parindent}{2em}
\setlength{\parskip}{0.28em}
\linespread{1.12}

\titleformat{\section}{\Large\bfseries\color{DeepTeal}}{}{0pt}{}
\titlespacing*{\section}{0pt}{0pt}{6pt}
\titleformat{\subsection}{\large\bfseries\color{DeepTeal}}{}{0pt}{}
\titlespacing*{\subsection}{0pt}{5pt}{3pt}

\pagestyle{fancy}
\fancyhf{}
\lhead{\small\color{DeepTeal}\textbf{文献阅读汇报}\quad <论文主题>}
\rhead{\small\color{gray}<作者>, <年>}
\cfoot{\small\color{gray}\thepage\,/\,\pageref{LastPage}}
\renewcommand{\headrulewidth}{0.5pt}
\renewcommand{\headrule}{\hbox to\headwidth{\color{DeepTeal}\leaders\hrule height \headrulewidth\hfill}}

% ---------- 版式组件 ----------
\newcommand{\key}[1]{\textcolor{DeepTeal}{\textbf{#1}}}
\newcommand{\highlight}[1]{\textcolor{WarmOrange}{\textbf{#1}}}
\newcommand{\source}[1]{\par\vspace{1pt}\raggedright\scriptsize\color{gray}\textit{来源：#1}\par}
\newcommand{\sectiontag}[1]{\noindent\colorbox{PaleTeal}{\parbox{\dimexpr\linewidth-2\fboxsep\relax}{\small\color{DeepTeal}\textbf{#1}}}\par\vspace{5pt}}

% 本节路线（浅青绿框）
\newcommand{\roadmap}[1]{\begin{tcolorbox}[colback=PaleTeal,colframe=MidTeal,boxrule=.5pt,arc=1.5mm,left=2.2mm,right=2.2mm,top=1.4mm,bottom=1.4mm]\small\textbf{本节路线：}\quad #1\end{tcolorbox}}
% 大白话（浅蓝框）
\newcommand{\plain}[1]{\begin{tcolorbox}[colback=PaleBlue,colframe=DeepTeal,boxrule=.5pt,arc=1.5mm,left=2.2mm,right=2.2mm,top=1.5mm,bottom=1.5mm]\small\textbf{大白话}\quad #1\end{tcolorbox}}
% 专业定义 + 本文中的位置（白底青绿框）
\newcommand{\definition}[2]{\begin{tcolorbox}[colback=white,colframe=MidTeal,boxrule=.45pt,arc=1.5mm,left=2.2mm,right=2.2mm,top=1.4mm,bottom=1.4mm]\small\textbf{专业定义：}#1\par\vspace{1mm}\textbf{本文中的位置：}#2\end{tcolorbox}}
% 不要这样理解（暖橙框）
\newcommand{\warning}[1]{\begin{tcolorbox}[colback=WarmOrange!10,colframe=WarmOrange,boxrule=.5pt,arc=1.5mm,left=2.2mm,right=2.2mm,top=1.4mm,bottom=1.4mm]\small\textbf{不要这样理解：}\quad #1\end{tcolorbox}}
% 阅读要点（青绿框）
\newcommand{\takeaway}[1]{\begin{tcolorbox}[colback=PaleTeal,colframe=MidTeal,boxrule=.55pt,arc=1.5mm,left=2.2mm,right=2.2mm,top=1.5mm,bottom=1.5mm]\small\textbf{阅读要点}\quad #1\end{tcolorbox}}
% 注意（白底暖橙框）
\newcommand{\noteBox}[1]{\begin{tcolorbox}[colback=white,colframe=WarmOrange,boxrule=.55pt,arc=1.5mm,left=2.2mm,right=2.2mm,top=1.3mm,bottom=1.3mm]\small #1\end{tcolorbox}}
% 论文图（紧跟图注）
\newcommand{\paperfig}[3]{\begin{figure}[H]\centering\includegraphics[width=#2\linewidth]{#1}\caption{#3}\end{figure}}
% 代码文件标注
\newcommand{\codefile}[2]{\noindent\scriptsize\color{gray}\texttt{#1}\quad\textit{#2}\par\vspace{2pt}}

% ---------- 代码样式 ----------
\lstdefinestyle{pythonstyle}{language=Python,basicstyle=\ttfamily\scriptsize\color{Ink},
  keywordstyle=\color{DeepTeal}\bfseries,commentstyle=\color{gray}\itshape,
  stringstyle=\color{WarmOrange},backgroundcolor=\color{CodeBg},frame=single,
  rulecolor=\color{CodeRule},framerule=.35pt,breaklines=true,showstringspaces=false,
  columns=fullflexible,aboveskip=3pt,belowskip=3pt,xleftmargin=1pt,xrightmargin=1pt}

% ---------- 表格列类型 ----------
\newcolumntype{Y}{>{\RaggedRight\arraybackslash}X}
\newcolumntype{C}{>{\Centering\arraybackslash}X}

% ---------- 安全表格（避免 p{}+Y 混用溢出；见 SKILL.md 陷阱 6） ----------
% 用法：
%   \begin{center}\small
%   \begin{tabular}{@{}p{0.26\linewidth}@{\hspace{4mm}}p{0.62\linewidth}@{}}
%   ...\end{tabular}
%   \end{center}

% 注意：不要用 \accent 作命令名（与 LaTeX 内置冲突），用 \highlight。
% ============================================================
\documentclass[11pt,a4paper]{ctexart}

\usepackage[a4paper,margin=19mm,headheight=23pt,headsep=7mm,footskip=10mm]{geometry}
\usepackage{fontspec}
\setmainfont{Times New Roman}
\setsansfont{Microsoft YaHei}
\setmonofont{Consolas}[Scale=0.84]
\usepackage{amsmath,amssymb}
\usepackage{graphicx}
\usepackage{booktabs,tabularx,array,multirow}
\usepackage{caption,subcaption}
\usepackage{longtable}
\usepackage{makecell}
\usepackage{enumitem}
\usepackage{float}
\usepackage{listings}
\usepackage{tcolorbox}
\tcbuselibrary{skins,breakable}
\usepackage{tikz}
\usetikzlibrary{arrows.meta,positioning,shapes.geometric,fit,backgrounds}
\usepackage{fancyhdr}
\usepackage{lastpage}
\usepackage{hyperref}
\usepackage{xcolor}
\usepackage{titlesec}
\usepackage{setspace}
\usepackage{ragged2e}
\usepackage{pifont}

\definecolor{DeepTeal}{HTML}{0E5A67}
\definecolor{MidTeal}{HTML}{3E8C88}
\definecolor{WarmOrange}{HTML}{D17A35}
\definecolor{Ink}{HTML}{1D2B30}
\definecolor{PaleBlue}{HTML}{F3F7F8}
\definecolor{PaleTeal}{HTML}{E7F2F1}
\definecolor{CodeBg}{HTML}{F5F7F7}
\definecolor{CodeRule}{HTML}{C6D5D6}

% 建议在正文 % ============================================================
% 论文带读 PDF 统一主题模板（XeLaTeX + ctexart）
% 使用：\input{preamble}
% 注意：不要用 \accent 作命令名（与 LaTeX 内置冲突），用 \highlight。
% ============================================================
\documentclass[11pt,a4paper]{ctexart}

\usepackage[a4paper,margin=19mm,headheight=23pt,headsep=7mm,footskip=10mm]{geometry}
\usepackage{fontspec}
\setmainfont{Times New Roman}
\setsansfont{Microsoft YaHei}
\setmonofont{Consolas}[Scale=0.84]
\usepackage{amsmath,amssymb}
\usepackage{graphicx}
\usepackage{booktabs,tabularx,array,multirow}
\usepackage{caption,subcaption}
\usepackage{longtable}
\usepackage{makecell}
\usepackage{enumitem}
\usepackage{float}
\usepackage{listings}
\usepackage{tcolorbox}
\tcbuselibrary{skins,breakable}
\usepackage{tikz}
\usetikzlibrary{arrows.meta,positioning,shapes.geometric,fit,backgrounds}
\usepackage{fancyhdr}
\usepackage{lastpage}
\usepackage{hyperref}
\usepackage{xcolor}
\usepackage{titlesec}
\usepackage{setspace}
\usepackage{ragged2e}
\usepackage{pifont}

\definecolor{DeepTeal}{HTML}{0E5A67}
\definecolor{MidTeal}{HTML}{3E8C88}
\definecolor{WarmOrange}{HTML}{D17A35}
\definecolor{Ink}{HTML}{1D2B30}
\definecolor{PaleBlue}{HTML}{F3F7F8}
\definecolor{PaleTeal}{HTML}{E7F2F1}
\definecolor{CodeBg}{HTML}{F5F7F7}
\definecolor{CodeRule}{HTML}{C6D5D6}

% 建议在正文 \input{preamble} 之后用 \hypersetup 覆盖标题
\hypersetup{colorlinks=true,linkcolor=DeepTeal,urlcolor=DeepTeal,citecolor=DeepTeal,
  pdftitle={论文带读},pdfauthor={论文带读}}
\captionsetup{font=small,labelfont={bf,color=DeepTeal},textfont={color=Ink},skip=5pt}
\setlength{\parindent}{2em}
\setlength{\parskip}{0.28em}
\linespread{1.12}

\titleformat{\section}{\Large\bfseries\color{DeepTeal}}{}{0pt}{}
\titlespacing*{\section}{0pt}{0pt}{6pt}
\titleformat{\subsection}{\large\bfseries\color{DeepTeal}}{}{0pt}{}
\titlespacing*{\subsection}{0pt}{5pt}{3pt}

\pagestyle{fancy}
\fancyhf{}
\lhead{\small\color{DeepTeal}\textbf{文献阅读汇报}\quad <论文主题>}
\rhead{\small\color{gray}<作者>, <年>}
\cfoot{\small\color{gray}\thepage\,/\,\pageref{LastPage}}
\renewcommand{\headrulewidth}{0.5pt}
\renewcommand{\headrule}{\hbox to\headwidth{\color{DeepTeal}\leaders\hrule height \headrulewidth\hfill}}

% ---------- 版式组件 ----------
\newcommand{\key}[1]{\textcolor{DeepTeal}{\textbf{#1}}}
\newcommand{\highlight}[1]{\textcolor{WarmOrange}{\textbf{#1}}}
\newcommand{\source}[1]{\par\vspace{1pt}\raggedright\scriptsize\color{gray}\textit{来源：#1}\par}
\newcommand{\sectiontag}[1]{\noindent\colorbox{PaleTeal}{\parbox{\dimexpr\linewidth-2\fboxsep\relax}{\small\color{DeepTeal}\textbf{#1}}}\par\vspace{5pt}}

% 本节路线（浅青绿框）
\newcommand{\roadmap}[1]{\begin{tcolorbox}[colback=PaleTeal,colframe=MidTeal,boxrule=.5pt,arc=1.5mm,left=2.2mm,right=2.2mm,top=1.4mm,bottom=1.4mm]\small\textbf{本节路线：}\quad #1\end{tcolorbox}}
% 大白话（浅蓝框）
\newcommand{\plain}[1]{\begin{tcolorbox}[colback=PaleBlue,colframe=DeepTeal,boxrule=.5pt,arc=1.5mm,left=2.2mm,right=2.2mm,top=1.5mm,bottom=1.5mm]\small\textbf{大白话}\quad #1\end{tcolorbox}}
% 专业定义 + 本文中的位置（白底青绿框）
\newcommand{\definition}[2]{\begin{tcolorbox}[colback=white,colframe=MidTeal,boxrule=.45pt,arc=1.5mm,left=2.2mm,right=2.2mm,top=1.4mm,bottom=1.4mm]\small\textbf{专业定义：}#1\par\vspace{1mm}\textbf{本文中的位置：}#2\end{tcolorbox}}
% 不要这样理解（暖橙框）
\newcommand{\warning}[1]{\begin{tcolorbox}[colback=WarmOrange!10,colframe=WarmOrange,boxrule=.5pt,arc=1.5mm,left=2.2mm,right=2.2mm,top=1.4mm,bottom=1.4mm]\small\textbf{不要这样理解：}\quad #1\end{tcolorbox}}
% 阅读要点（青绿框）
\newcommand{\takeaway}[1]{\begin{tcolorbox}[colback=PaleTeal,colframe=MidTeal,boxrule=.55pt,arc=1.5mm,left=2.2mm,right=2.2mm,top=1.5mm,bottom=1.5mm]\small\textbf{阅读要点}\quad #1\end{tcolorbox}}
% 注意（白底暖橙框）
\newcommand{\noteBox}[1]{\begin{tcolorbox}[colback=white,colframe=WarmOrange,boxrule=.55pt,arc=1.5mm,left=2.2mm,right=2.2mm,top=1.3mm,bottom=1.3mm]\small #1\end{tcolorbox}}
% 论文图（紧跟图注）
\newcommand{\paperfig}[3]{\begin{figure}[H]\centering\includegraphics[width=#2\linewidth]{#1}\caption{#3}\end{figure}}
% 代码文件标注
\newcommand{\codefile}[2]{\noindent\scriptsize\color{gray}\texttt{#1}\quad\textit{#2}\par\vspace{2pt}}

% ---------- 代码样式 ----------
\lstdefinestyle{pythonstyle}{language=Python,basicstyle=\ttfamily\scriptsize\color{Ink},
  keywordstyle=\color{DeepTeal}\bfseries,commentstyle=\color{gray}\itshape,
  stringstyle=\color{WarmOrange},backgroundcolor=\color{CodeBg},frame=single,
  rulecolor=\color{CodeRule},framerule=.35pt,breaklines=true,showstringspaces=false,
  columns=fullflexible,aboveskip=3pt,belowskip=3pt,xleftmargin=1pt,xrightmargin=1pt}

% ---------- 表格列类型 ----------
\newcolumntype{Y}{>{\RaggedRight\arraybackslash}X}
\newcolumntype{C}{>{\Centering\arraybackslash}X}

% ---------- 安全表格（避免 p{}+Y 混用溢出；见 SKILL.md 陷阱 6） ----------
% 用法：
%   \begin{center}\small
%   \begin{tabular}{@{}p{0.26\linewidth}@{\hspace{4mm}}p{0.62\linewidth}@{}}
%   ...\end{tabular}
%   \end{center}
 之后用 \hypersetup 覆盖标题
\hypersetup{colorlinks=true,linkcolor=DeepTeal,urlcolor=DeepTeal,citecolor=DeepTeal,
  pdftitle={论文带读},pdfauthor={论文带读}}
\captionsetup{font=small,labelfont={bf,color=DeepTeal},textfont={color=Ink},skip=5pt}
\setlength{\parindent}{2em}
\setlength{\parskip}{0.28em}
\linespread{1.12}

\titleformat{\section}{\Large\bfseries\color{DeepTeal}}{}{0pt}{}
\titlespacing*{\section}{0pt}{0pt}{6pt}
\titleformat{\subsection}{\large\bfseries\color{DeepTeal}}{}{0pt}{}
\titlespacing*{\subsection}{0pt}{5pt}{3pt}

\pagestyle{fancy}
\fancyhf{}
\lhead{\small\color{DeepTeal}\textbf{文献阅读汇报}\quad <论文主题>}
\rhead{\small\color{gray}<作者>, <年>}
\cfoot{\small\color{gray}\thepage\,/\,\pageref{LastPage}}
\renewcommand{\headrulewidth}{0.5pt}
\renewcommand{\headrule}{\hbox to\headwidth{\color{DeepTeal}\leaders\hrule height \headrulewidth\hfill}}

% ---------- 版式组件 ----------
\newcommand{\key}[1]{\textcolor{DeepTeal}{\textbf{#1}}}
\newcommand{\highlight}[1]{\textcolor{WarmOrange}{\textbf{#1}}}
\newcommand{\source}[1]{\par\vspace{1pt}\raggedright\scriptsize\color{gray}\textit{来源：#1}\par}
\newcommand{\sectiontag}[1]{\noindent\colorbox{PaleTeal}{\parbox{\dimexpr\linewidth-2\fboxsep\relax}{\small\color{DeepTeal}\textbf{#1}}}\par\vspace{5pt}}

% 本节路线（浅青绿框）
\newcommand{\roadmap}[1]{\begin{tcolorbox}[colback=PaleTeal,colframe=MidTeal,boxrule=.5pt,arc=1.5mm,left=2.2mm,right=2.2mm,top=1.4mm,bottom=1.4mm]\small\textbf{本节路线：}\quad #1\end{tcolorbox}}
% 大白话（浅蓝框）
\newcommand{\plain}[1]{\begin{tcolorbox}[colback=PaleBlue,colframe=DeepTeal,boxrule=.5pt,arc=1.5mm,left=2.2mm,right=2.2mm,top=1.5mm,bottom=1.5mm]\small\textbf{大白话}\quad #1\end{tcolorbox}}
% 专业定义 + 本文中的位置（白底青绿框）
\newcommand{\definition}[2]{\begin{tcolorbox}[colback=white,colframe=MidTeal,boxrule=.45pt,arc=1.5mm,left=2.2mm,right=2.2mm,top=1.4mm,bottom=1.4mm]\small\textbf{专业定义：}#1\par\vspace{1mm}\textbf{本文中的位置：}#2\end{tcolorbox}}
% 不要这样理解（暖橙框）
\newcommand{\warning}[1]{\begin{tcolorbox}[colback=WarmOrange!10,colframe=WarmOrange,boxrule=.5pt,arc=1.5mm,left=2.2mm,right=2.2mm,top=1.4mm,bottom=1.4mm]\small\textbf{不要这样理解：}\quad #1\end{tcolorbox}}
% 阅读要点（青绿框）
\newcommand{\takeaway}[1]{\begin{tcolorbox}[colback=PaleTeal,colframe=MidTeal,boxrule=.55pt,arc=1.5mm,left=2.2mm,right=2.2mm,top=1.5mm,bottom=1.5mm]\small\textbf{阅读要点}\quad #1\end{tcolorbox}}
% 注意（白底暖橙框）
\newcommand{\noteBox}[1]{\begin{tcolorbox}[colback=white,colframe=WarmOrange,boxrule=.55pt,arc=1.5mm,left=2.2mm,right=2.2mm,top=1.3mm,bottom=1.3mm]\small #1\end{tcolorbox}}
% 论文图（紧跟图注）
\newcommand{\paperfig}[3]{\begin{figure}[H]\centering\includegraphics[width=#2\linewidth]{#1}\caption{#3}\end{figure}}
% 代码文件标注
\newcommand{\codefile}[2]{\noindent\scriptsize\color{gray}\texttt{#1}\quad\textit{#2}\par\vspace{2pt}}

% ---------- 代码样式 ----------
\lstdefinestyle{pythonstyle}{language=Python,basicstyle=\ttfamily\scriptsize\color{Ink},
  keywordstyle=\color{DeepTeal}\bfseries,commentstyle=\color{gray}\itshape,
  stringstyle=\color{WarmOrange},backgroundcolor=\color{CodeBg},frame=single,
  rulecolor=\color{CodeRule},framerule=.35pt,breaklines=true,showstringspaces=false,
  columns=fullflexible,aboveskip=3pt,belowskip=3pt,xleftmargin=1pt,xrightmargin=1pt}

% ---------- 表格列类型 ----------
\newcolumntype{Y}{>{\RaggedRight\arraybackslash}X}
\newcolumntype{C}{>{\Centering\arraybackslash}X}

% ---------- 安全表格（避免 p{}+Y 混用溢出；见 SKILL.md 陷阱 6） ----------
% 用法：
%   \begin{center}\small
%   \begin{tabular}{@{}p{0.26\linewidth}@{\hspace{4mm}}p{0.62\linewidth}@{}}
%   ...\end{tabular}
%   \end{center}

% 注意：不要用 \accent 作命令名（与 LaTeX 内置冲突），用 \highlight。
% ============================================================
\documentclass[11pt,a4paper]{ctexart}

\usepackage[a4paper,margin=19mm,headheight=23pt,headsep=7mm,footskip=10mm]{geometry}
\usepackage{fontspec}
\setmainfont{Times New Roman}
\setsansfont{Microsoft YaHei}
\setmonofont{Consolas}[Scale=0.84]
\usepackage{amsmath,amssymb}
\usepackage{graphicx}
\usepackage{booktabs,tabularx,array,multirow}
\usepackage{caption,subcaption}
\usepackage{longtable}
\usepackage{makecell}
\usepackage{enumitem}
\usepackage{float}
\usepackage{listings}
\usepackage{tcolorbox}
\tcbuselibrary{skins,breakable}
\usepackage{tikz}
\usetikzlibrary{arrows.meta,positioning,shapes.geometric,fit,backgrounds}
\usepackage{fancyhdr}
\usepackage{lastpage}
\usepackage{hyperref}
\usepackage{xcolor}
\usepackage{titlesec}
\usepackage{setspace}
\usepackage{ragged2e}
\usepackage{pifont}

\definecolor{DeepTeal}{HTML}{0E5A67}
\definecolor{MidTeal}{HTML}{3E8C88}
\definecolor{WarmOrange}{HTML}{D17A35}
\definecolor{Ink}{HTML}{1D2B30}
\definecolor{PaleBlue}{HTML}{F3F7F8}
\definecolor{PaleTeal}{HTML}{E7F2F1}
\definecolor{CodeBg}{HTML}{F5F7F7}
\definecolor{CodeRule}{HTML}{C6D5D6}

% 建议在正文 % ============================================================
% 论文带读 PDF 统一主题模板（XeLaTeX + ctexart）
% 使用：% ============================================================
% 论文带读 PDF 统一主题模板（XeLaTeX + ctexart）
% 使用：\input{preamble}
% 注意：不要用 \accent 作命令名（与 LaTeX 内置冲突），用 \highlight。
% ============================================================
\documentclass[11pt,a4paper]{ctexart}

\usepackage[a4paper,margin=19mm,headheight=23pt,headsep=7mm,footskip=10mm]{geometry}
\usepackage{fontspec}
\setmainfont{Times New Roman}
\setsansfont{Microsoft YaHei}
\setmonofont{Consolas}[Scale=0.84]
\usepackage{amsmath,amssymb}
\usepackage{graphicx}
\usepackage{booktabs,tabularx,array,multirow}
\usepackage{caption,subcaption}
\usepackage{longtable}
\usepackage{makecell}
\usepackage{enumitem}
\usepackage{float}
\usepackage{listings}
\usepackage{tcolorbox}
\tcbuselibrary{skins,breakable}
\usepackage{tikz}
\usetikzlibrary{arrows.meta,positioning,shapes.geometric,fit,backgrounds}
\usepackage{fancyhdr}
\usepackage{lastpage}
\usepackage{hyperref}
\usepackage{xcolor}
\usepackage{titlesec}
\usepackage{setspace}
\usepackage{ragged2e}
\usepackage{pifont}

\definecolor{DeepTeal}{HTML}{0E5A67}
\definecolor{MidTeal}{HTML}{3E8C88}
\definecolor{WarmOrange}{HTML}{D17A35}
\definecolor{Ink}{HTML}{1D2B30}
\definecolor{PaleBlue}{HTML}{F3F7F8}
\definecolor{PaleTeal}{HTML}{E7F2F1}
\definecolor{CodeBg}{HTML}{F5F7F7}
\definecolor{CodeRule}{HTML}{C6D5D6}

% 建议在正文 \input{preamble} 之后用 \hypersetup 覆盖标题
\hypersetup{colorlinks=true,linkcolor=DeepTeal,urlcolor=DeepTeal,citecolor=DeepTeal,
  pdftitle={论文带读},pdfauthor={论文带读}}
\captionsetup{font=small,labelfont={bf,color=DeepTeal},textfont={color=Ink},skip=5pt}
\setlength{\parindent}{2em}
\setlength{\parskip}{0.28em}
\linespread{1.12}

\titleformat{\section}{\Large\bfseries\color{DeepTeal}}{}{0pt}{}
\titlespacing*{\section}{0pt}{0pt}{6pt}
\titleformat{\subsection}{\large\bfseries\color{DeepTeal}}{}{0pt}{}
\titlespacing*{\subsection}{0pt}{5pt}{3pt}

\pagestyle{fancy}
\fancyhf{}
\lhead{\small\color{DeepTeal}\textbf{文献阅读汇报}\quad <论文主题>}
\rhead{\small\color{gray}<作者>, <年>}
\cfoot{\small\color{gray}\thepage\,/\,\pageref{LastPage}}
\renewcommand{\headrulewidth}{0.5pt}
\renewcommand{\headrule}{\hbox to\headwidth{\color{DeepTeal}\leaders\hrule height \headrulewidth\hfill}}

% ---------- 版式组件 ----------
\newcommand{\key}[1]{\textcolor{DeepTeal}{\textbf{#1}}}
\newcommand{\highlight}[1]{\textcolor{WarmOrange}{\textbf{#1}}}
\newcommand{\source}[1]{\par\vspace{1pt}\raggedright\scriptsize\color{gray}\textit{来源：#1}\par}
\newcommand{\sectiontag}[1]{\noindent\colorbox{PaleTeal}{\parbox{\dimexpr\linewidth-2\fboxsep\relax}{\small\color{DeepTeal}\textbf{#1}}}\par\vspace{5pt}}

% 本节路线（浅青绿框）
\newcommand{\roadmap}[1]{\begin{tcolorbox}[colback=PaleTeal,colframe=MidTeal,boxrule=.5pt,arc=1.5mm,left=2.2mm,right=2.2mm,top=1.4mm,bottom=1.4mm]\small\textbf{本节路线：}\quad #1\end{tcolorbox}}
% 大白话（浅蓝框）
\newcommand{\plain}[1]{\begin{tcolorbox}[colback=PaleBlue,colframe=DeepTeal,boxrule=.5pt,arc=1.5mm,left=2.2mm,right=2.2mm,top=1.5mm,bottom=1.5mm]\small\textbf{大白话}\quad #1\end{tcolorbox}}
% 专业定义 + 本文中的位置（白底青绿框）
\newcommand{\definition}[2]{\begin{tcolorbox}[colback=white,colframe=MidTeal,boxrule=.45pt,arc=1.5mm,left=2.2mm,right=2.2mm,top=1.4mm,bottom=1.4mm]\small\textbf{专业定义：}#1\par\vspace{1mm}\textbf{本文中的位置：}#2\end{tcolorbox}}
% 不要这样理解（暖橙框）
\newcommand{\warning}[1]{\begin{tcolorbox}[colback=WarmOrange!10,colframe=WarmOrange,boxrule=.5pt,arc=1.5mm,left=2.2mm,right=2.2mm,top=1.4mm,bottom=1.4mm]\small\textbf{不要这样理解：}\quad #1\end{tcolorbox}}
% 阅读要点（青绿框）
\newcommand{\takeaway}[1]{\begin{tcolorbox}[colback=PaleTeal,colframe=MidTeal,boxrule=.55pt,arc=1.5mm,left=2.2mm,right=2.2mm,top=1.5mm,bottom=1.5mm]\small\textbf{阅读要点}\quad #1\end{tcolorbox}}
% 注意（白底暖橙框）
\newcommand{\noteBox}[1]{\begin{tcolorbox}[colback=white,colframe=WarmOrange,boxrule=.55pt,arc=1.5mm,left=2.2mm,right=2.2mm,top=1.3mm,bottom=1.3mm]\small #1\end{tcolorbox}}
% 论文图（紧跟图注）
\newcommand{\paperfig}[3]{\begin{figure}[H]\centering\includegraphics[width=#2\linewidth]{#1}\caption{#3}\end{figure}}
% 代码文件标注
\newcommand{\codefile}[2]{\noindent\scriptsize\color{gray}\texttt{#1}\quad\textit{#2}\par\vspace{2pt}}

% ---------- 代码样式 ----------
\lstdefinestyle{pythonstyle}{language=Python,basicstyle=\ttfamily\scriptsize\color{Ink},
  keywordstyle=\color{DeepTeal}\bfseries,commentstyle=\color{gray}\itshape,
  stringstyle=\color{WarmOrange},backgroundcolor=\color{CodeBg},frame=single,
  rulecolor=\color{CodeRule},framerule=.35pt,breaklines=true,showstringspaces=false,
  columns=fullflexible,aboveskip=3pt,belowskip=3pt,xleftmargin=1pt,xrightmargin=1pt}

% ---------- 表格列类型 ----------
\newcolumntype{Y}{>{\RaggedRight\arraybackslash}X}
\newcolumntype{C}{>{\Centering\arraybackslash}X}

% ---------- 安全表格（避免 p{}+Y 混用溢出；见 SKILL.md 陷阱 6） ----------
% 用法：
%   \begin{center}\small
%   \begin{tabular}{@{}p{0.26\linewidth}@{\hspace{4mm}}p{0.62\linewidth}@{}}
%   ...\end{tabular}
%   \end{center}

% 注意：不要用 \accent 作命令名（与 LaTeX 内置冲突），用 \highlight。
% ============================================================
\documentclass[11pt,a4paper]{ctexart}

\usepackage[a4paper,margin=19mm,headheight=23pt,headsep=7mm,footskip=10mm]{geometry}
\usepackage{fontspec}
\setmainfont{Times New Roman}
\setsansfont{Microsoft YaHei}
\setmonofont{Consolas}[Scale=0.84]
\usepackage{amsmath,amssymb}
\usepackage{graphicx}
\usepackage{booktabs,tabularx,array,multirow}
\usepackage{caption,subcaption}
\usepackage{longtable}
\usepackage{makecell}
\usepackage{enumitem}
\usepackage{float}
\usepackage{listings}
\usepackage{tcolorbox}
\tcbuselibrary{skins,breakable}
\usepackage{tikz}
\usetikzlibrary{arrows.meta,positioning,shapes.geometric,fit,backgrounds}
\usepackage{fancyhdr}
\usepackage{lastpage}
\usepackage{hyperref}
\usepackage{xcolor}
\usepackage{titlesec}
\usepackage{setspace}
\usepackage{ragged2e}
\usepackage{pifont}

\definecolor{DeepTeal}{HTML}{0E5A67}
\definecolor{MidTeal}{HTML}{3E8C88}
\definecolor{WarmOrange}{HTML}{D17A35}
\definecolor{Ink}{HTML}{1D2B30}
\definecolor{PaleBlue}{HTML}{F3F7F8}
\definecolor{PaleTeal}{HTML}{E7F2F1}
\definecolor{CodeBg}{HTML}{F5F7F7}
\definecolor{CodeRule}{HTML}{C6D5D6}

% 建议在正文 % ============================================================
% 论文带读 PDF 统一主题模板（XeLaTeX + ctexart）
% 使用：\input{preamble}
% 注意：不要用 \accent 作命令名（与 LaTeX 内置冲突），用 \highlight。
% ============================================================
\documentclass[11pt,a4paper]{ctexart}

\usepackage[a4paper,margin=19mm,headheight=23pt,headsep=7mm,footskip=10mm]{geometry}
\usepackage{fontspec}
\setmainfont{Times New Roman}
\setsansfont{Microsoft YaHei}
\setmonofont{Consolas}[Scale=0.84]
\usepackage{amsmath,amssymb}
\usepackage{graphicx}
\usepackage{booktabs,tabularx,array,multirow}
\usepackage{caption,subcaption}
\usepackage{longtable}
\usepackage{makecell}
\usepackage{enumitem}
\usepackage{float}
\usepackage{listings}
\usepackage{tcolorbox}
\tcbuselibrary{skins,breakable}
\usepackage{tikz}
\usetikzlibrary{arrows.meta,positioning,shapes.geometric,fit,backgrounds}
\usepackage{fancyhdr}
\usepackage{lastpage}
\usepackage{hyperref}
\usepackage{xcolor}
\usepackage{titlesec}
\usepackage{setspace}
\usepackage{ragged2e}
\usepackage{pifont}

\definecolor{DeepTeal}{HTML}{0E5A67}
\definecolor{MidTeal}{HTML}{3E8C88}
\definecolor{WarmOrange}{HTML}{D17A35}
\definecolor{Ink}{HTML}{1D2B30}
\definecolor{PaleBlue}{HTML}{F3F7F8}
\definecolor{PaleTeal}{HTML}{E7F2F1}
\definecolor{CodeBg}{HTML}{F5F7F7}
\definecolor{CodeRule}{HTML}{C6D5D6}

% 建议在正文 \input{preamble} 之后用 \hypersetup 覆盖标题
\hypersetup{colorlinks=true,linkcolor=DeepTeal,urlcolor=DeepTeal,citecolor=DeepTeal,
  pdftitle={论文带读},pdfauthor={论文带读}}
\captionsetup{font=small,labelfont={bf,color=DeepTeal},textfont={color=Ink},skip=5pt}
\setlength{\parindent}{2em}
\setlength{\parskip}{0.28em}
\linespread{1.12}

\titleformat{\section}{\Large\bfseries\color{DeepTeal}}{}{0pt}{}
\titlespacing*{\section}{0pt}{0pt}{6pt}
\titleformat{\subsection}{\large\bfseries\color{DeepTeal}}{}{0pt}{}
\titlespacing*{\subsection}{0pt}{5pt}{3pt}

\pagestyle{fancy}
\fancyhf{}
\lhead{\small\color{DeepTeal}\textbf{文献阅读汇报}\quad <论文主题>}
\rhead{\small\color{gray}<作者>, <年>}
\cfoot{\small\color{gray}\thepage\,/\,\pageref{LastPage}}
\renewcommand{\headrulewidth}{0.5pt}
\renewcommand{\headrule}{\hbox to\headwidth{\color{DeepTeal}\leaders\hrule height \headrulewidth\hfill}}

% ---------- 版式组件 ----------
\newcommand{\key}[1]{\textcolor{DeepTeal}{\textbf{#1}}}
\newcommand{\highlight}[1]{\textcolor{WarmOrange}{\textbf{#1}}}
\newcommand{\source}[1]{\par\vspace{1pt}\raggedright\scriptsize\color{gray}\textit{来源：#1}\par}
\newcommand{\sectiontag}[1]{\noindent\colorbox{PaleTeal}{\parbox{\dimexpr\linewidth-2\fboxsep\relax}{\small\color{DeepTeal}\textbf{#1}}}\par\vspace{5pt}}

% 本节路线（浅青绿框）
\newcommand{\roadmap}[1]{\begin{tcolorbox}[colback=PaleTeal,colframe=MidTeal,boxrule=.5pt,arc=1.5mm,left=2.2mm,right=2.2mm,top=1.4mm,bottom=1.4mm]\small\textbf{本节路线：}\quad #1\end{tcolorbox}}
% 大白话（浅蓝框）
\newcommand{\plain}[1]{\begin{tcolorbox}[colback=PaleBlue,colframe=DeepTeal,boxrule=.5pt,arc=1.5mm,left=2.2mm,right=2.2mm,top=1.5mm,bottom=1.5mm]\small\textbf{大白话}\quad #1\end{tcolorbox}}
% 专业定义 + 本文中的位置（白底青绿框）
\newcommand{\definition}[2]{\begin{tcolorbox}[colback=white,colframe=MidTeal,boxrule=.45pt,arc=1.5mm,left=2.2mm,right=2.2mm,top=1.4mm,bottom=1.4mm]\small\textbf{专业定义：}#1\par\vspace{1mm}\textbf{本文中的位置：}#2\end{tcolorbox}}
% 不要这样理解（暖橙框）
\newcommand{\warning}[1]{\begin{tcolorbox}[colback=WarmOrange!10,colframe=WarmOrange,boxrule=.5pt,arc=1.5mm,left=2.2mm,right=2.2mm,top=1.4mm,bottom=1.4mm]\small\textbf{不要这样理解：}\quad #1\end{tcolorbox}}
% 阅读要点（青绿框）
\newcommand{\takeaway}[1]{\begin{tcolorbox}[colback=PaleTeal,colframe=MidTeal,boxrule=.55pt,arc=1.5mm,left=2.2mm,right=2.2mm,top=1.5mm,bottom=1.5mm]\small\textbf{阅读要点}\quad #1\end{tcolorbox}}
% 注意（白底暖橙框）
\newcommand{\noteBox}[1]{\begin{tcolorbox}[colback=white,colframe=WarmOrange,boxrule=.55pt,arc=1.5mm,left=2.2mm,right=2.2mm,top=1.3mm,bottom=1.3mm]\small #1\end{tcolorbox}}
% 论文图（紧跟图注）
\newcommand{\paperfig}[3]{\begin{figure}[H]\centering\includegraphics[width=#2\linewidth]{#1}\caption{#3}\end{figure}}
% 代码文件标注
\newcommand{\codefile}[2]{\noindent\scriptsize\color{gray}\texttt{#1}\quad\textit{#2}\par\vspace{2pt}}

% ---------- 代码样式 ----------
\lstdefinestyle{pythonstyle}{language=Python,basicstyle=\ttfamily\scriptsize\color{Ink},
  keywordstyle=\color{DeepTeal}\bfseries,commentstyle=\color{gray}\itshape,
  stringstyle=\color{WarmOrange},backgroundcolor=\color{CodeBg},frame=single,
  rulecolor=\color{CodeRule},framerule=.35pt,breaklines=true,showstringspaces=false,
  columns=fullflexible,aboveskip=3pt,belowskip=3pt,xleftmargin=1pt,xrightmargin=1pt}

% ---------- 表格列类型 ----------
\newcolumntype{Y}{>{\RaggedRight\arraybackslash}X}
\newcolumntype{C}{>{\Centering\arraybackslash}X}

% ---------- 安全表格（避免 p{}+Y 混用溢出；见 SKILL.md 陷阱 6） ----------
% 用法：
%   \begin{center}\small
%   \begin{tabular}{@{}p{0.26\linewidth}@{\hspace{4mm}}p{0.62\linewidth}@{}}
%   ...\end{tabular}
%   \end{center}
 之后用 \hypersetup 覆盖标题
\hypersetup{colorlinks=true,linkcolor=DeepTeal,urlcolor=DeepTeal,citecolor=DeepTeal,
  pdftitle={论文带读},pdfauthor={论文带读}}
\captionsetup{font=small,labelfont={bf,color=DeepTeal},textfont={color=Ink},skip=5pt}
\setlength{\parindent}{2em}
\setlength{\parskip}{0.28em}
\linespread{1.12}

\titleformat{\section}{\Large\bfseries\color{DeepTeal}}{}{0pt}{}
\titlespacing*{\section}{0pt}{0pt}{6pt}
\titleformat{\subsection}{\large\bfseries\color{DeepTeal}}{}{0pt}{}
\titlespacing*{\subsection}{0pt}{5pt}{3pt}

\pagestyle{fancy}
\fancyhf{}
\lhead{\small\color{DeepTeal}\textbf{文献阅读汇报}\quad <论文主题>}
\rhead{\small\color{gray}<作者>, <年>}
\cfoot{\small\color{gray}\thepage\,/\,\pageref{LastPage}}
\renewcommand{\headrulewidth}{0.5pt}
\renewcommand{\headrule}{\hbox to\headwidth{\color{DeepTeal}\leaders\hrule height \headrulewidth\hfill}}

% ---------- 版式组件 ----------
\newcommand{\key}[1]{\textcolor{DeepTeal}{\textbf{#1}}}
\newcommand{\highlight}[1]{\textcolor{WarmOrange}{\textbf{#1}}}
\newcommand{\source}[1]{\par\vspace{1pt}\raggedright\scriptsize\color{gray}\textit{来源：#1}\par}
\newcommand{\sectiontag}[1]{\noindent\colorbox{PaleTeal}{\parbox{\dimexpr\linewidth-2\fboxsep\relax}{\small\color{DeepTeal}\textbf{#1}}}\par\vspace{5pt}}

% 本节路线（浅青绿框）
\newcommand{\roadmap}[1]{\begin{tcolorbox}[colback=PaleTeal,colframe=MidTeal,boxrule=.5pt,arc=1.5mm,left=2.2mm,right=2.2mm,top=1.4mm,bottom=1.4mm]\small\textbf{本节路线：}\quad #1\end{tcolorbox}}
% 大白话（浅蓝框）
\newcommand{\plain}[1]{\begin{tcolorbox}[colback=PaleBlue,colframe=DeepTeal,boxrule=.5pt,arc=1.5mm,left=2.2mm,right=2.2mm,top=1.5mm,bottom=1.5mm]\small\textbf{大白话}\quad #1\end{tcolorbox}}
% 专业定义 + 本文中的位置（白底青绿框）
\newcommand{\definition}[2]{\begin{tcolorbox}[colback=white,colframe=MidTeal,boxrule=.45pt,arc=1.5mm,left=2.2mm,right=2.2mm,top=1.4mm,bottom=1.4mm]\small\textbf{专业定义：}#1\par\vspace{1mm}\textbf{本文中的位置：}#2\end{tcolorbox}}
% 不要这样理解（暖橙框）
\newcommand{\warning}[1]{\begin{tcolorbox}[colback=WarmOrange!10,colframe=WarmOrange,boxrule=.5pt,arc=1.5mm,left=2.2mm,right=2.2mm,top=1.4mm,bottom=1.4mm]\small\textbf{不要这样理解：}\quad #1\end{tcolorbox}}
% 阅读要点（青绿框）
\newcommand{\takeaway}[1]{\begin{tcolorbox}[colback=PaleTeal,colframe=MidTeal,boxrule=.55pt,arc=1.5mm,left=2.2mm,right=2.2mm,top=1.5mm,bottom=1.5mm]\small\textbf{阅读要点}\quad #1\end{tcolorbox}}
% 注意（白底暖橙框）
\newcommand{\noteBox}[1]{\begin{tcolorbox}[colback=white,colframe=WarmOrange,boxrule=.55pt,arc=1.5mm,left=2.2mm,right=2.2mm,top=1.3mm,bottom=1.3mm]\small #1\end{tcolorbox}}
% 论文图（紧跟图注）
\newcommand{\paperfig}[3]{\begin{figure}[H]\centering\includegraphics[width=#2\linewidth]{#1}\caption{#3}\end{figure}}
% 代码文件标注
\newcommand{\codefile}[2]{\noindent\scriptsize\color{gray}\texttt{#1}\quad\textit{#2}\par\vspace{2pt}}

% ---------- 代码样式 ----------
\lstdefinestyle{pythonstyle}{language=Python,basicstyle=\ttfamily\scriptsize\color{Ink},
  keywordstyle=\color{DeepTeal}\bfseries,commentstyle=\color{gray}\itshape,
  stringstyle=\color{WarmOrange},backgroundcolor=\color{CodeBg},frame=single,
  rulecolor=\color{CodeRule},framerule=.35pt,breaklines=true,showstringspaces=false,
  columns=fullflexible,aboveskip=3pt,belowskip=3pt,xleftmargin=1pt,xrightmargin=1pt}

% ---------- 表格列类型 ----------
\newcolumntype{Y}{>{\RaggedRight\arraybackslash}X}
\newcolumntype{C}{>{\Centering\arraybackslash}X}

% ---------- 安全表格（避免 p{}+Y 混用溢出；见 SKILL.md 陷阱 6） ----------
% 用法：
%   \begin{center}\small
%   \begin{tabular}{@{}p{0.26\linewidth}@{\hspace{4mm}}p{0.62\linewidth}@{}}
%   ...\end{tabular}
%   \end{center}
 之后用 \hypersetup 覆盖标题
\hypersetup{colorlinks=true,linkcolor=DeepTeal,urlcolor=DeepTeal,citecolor=DeepTeal,
  pdftitle={论文带读},pdfauthor={论文带读}}
\captionsetup{font=small,labelfont={bf,color=DeepTeal},textfont={color=Ink},skip=5pt}
\setlength{\parindent}{2em}
\setlength{\parskip}{0.28em}
\linespread{1.12}

\titleformat{\section}{\Large\bfseries\color{DeepTeal}}{}{0pt}{}
\titlespacing*{\section}{0pt}{0pt}{6pt}
\titleformat{\subsection}{\large\bfseries\color{DeepTeal}}{}{0pt}{}
\titlespacing*{\subsection}{0pt}{5pt}{3pt}

\pagestyle{fancy}
\fancyhf{}
\lhead{\small\color{DeepTeal}\textbf{文献阅读汇报}\quad <论文主题>}
\rhead{\small\color{gray}<作者>, <年>}
\cfoot{\small\color{gray}\thepage\,/\,\pageref{LastPage}}
\renewcommand{\headrulewidth}{0.5pt}
\renewcommand{\headrule}{\hbox to\headwidth{\color{DeepTeal}\leaders\hrule height \headrulewidth\hfill}}

% ---------- 版式组件 ----------
\newcommand{\key}[1]{\textcolor{DeepTeal}{\textbf{#1}}}
\newcommand{\highlight}[1]{\textcolor{WarmOrange}{\textbf{#1}}}
\newcommand{\source}[1]{\par\vspace{1pt}\raggedright\scriptsize\color{gray}\textit{来源：#1}\par}
\newcommand{\sectiontag}[1]{\noindent\colorbox{PaleTeal}{\parbox{\dimexpr\linewidth-2\fboxsep\relax}{\small\color{DeepTeal}\textbf{#1}}}\par\vspace{5pt}}

% 本节路线（浅青绿框）
\newcommand{\roadmap}[1]{\begin{tcolorbox}[colback=PaleTeal,colframe=MidTeal,boxrule=.5pt,arc=1.5mm,left=2.2mm,right=2.2mm,top=1.4mm,bottom=1.4mm]\small\textbf{本节路线：}\quad #1\end{tcolorbox}}
% 大白话（浅蓝框）
\newcommand{\plain}[1]{\begin{tcolorbox}[colback=PaleBlue,colframe=DeepTeal,boxrule=.5pt,arc=1.5mm,left=2.2mm,right=2.2mm,top=1.5mm,bottom=1.5mm]\small\textbf{大白话}\quad #1\end{tcolorbox}}
% 专业定义 + 本文中的位置（白底青绿框）
\newcommand{\definition}[2]{\begin{tcolorbox}[colback=white,colframe=MidTeal,boxrule=.45pt,arc=1.5mm,left=2.2mm,right=2.2mm,top=1.4mm,bottom=1.4mm]\small\textbf{专业定义：}#1\par\vspace{1mm}\textbf{本文中的位置：}#2\end{tcolorbox}}
% 不要这样理解（暖橙框）
\newcommand{\warning}[1]{\begin{tcolorbox}[colback=WarmOrange!10,colframe=WarmOrange,boxrule=.5pt,arc=1.5mm,left=2.2mm,right=2.2mm,top=1.4mm,bottom=1.4mm]\small\textbf{不要这样理解：}\quad #1\end{tcolorbox}}
% 阅读要点（青绿框）
\newcommand{\takeaway}[1]{\begin{tcolorbox}[colback=PaleTeal,colframe=MidTeal,boxrule=.55pt,arc=1.5mm,left=2.2mm,right=2.2mm,top=1.5mm,bottom=1.5mm]\small\textbf{阅读要点}\quad #1\end{tcolorbox}}
% 注意（白底暖橙框）
\newcommand{\noteBox}[1]{\begin{tcolorbox}[colback=white,colframe=WarmOrange,boxrule=.55pt,arc=1.5mm,left=2.2mm,right=2.2mm,top=1.3mm,bottom=1.3mm]\small #1\end{tcolorbox}}
% 论文图（紧跟图注）
\newcommand{\paperfig}[3]{\begin{figure}[H]\centering\includegraphics[width=#2\linewidth]{#1}\caption{#3}\end{figure}}
% 代码文件标注
\newcommand{\codefile}[2]{\noindent\scriptsize\color{gray}\texttt{#1}\quad\textit{#2}\par\vspace{2pt}}

% ---------- 代码样式 ----------
\lstdefinestyle{pythonstyle}{language=Python,basicstyle=\ttfamily\scriptsize\color{Ink},
  keywordstyle=\color{DeepTeal}\bfseries,commentstyle=\color{gray}\itshape,
  stringstyle=\color{WarmOrange},backgroundcolor=\color{CodeBg},frame=single,
  rulecolor=\color{CodeRule},framerule=.35pt,breaklines=true,showstringspaces=false,
  columns=fullflexible,aboveskip=3pt,belowskip=3pt,xleftmargin=1pt,xrightmargin=1pt}

% ---------- 表格列类型 ----------
\newcolumntype{Y}{>{\RaggedRight\arraybackslash}X}
\newcolumntype{C}{>{\Centering\arraybackslash}X}

% ---------- 安全表格（避免 p{}+Y 混用溢出；见 SKILL.md 陷阱 6） ----------
% 用法：
%   \begin{center}\small
%   \begin{tabular}{@{}p{0.26\linewidth}@{\hspace{4mm}}p{0.62\linewidth}@{}}
%   ...\end{tabular}
%   \end{center}
 之后用 \hypersetup 覆盖标题
\hypersetup{colorlinks=true,linkcolor=DeepTeal,urlcolor=DeepTeal,citecolor=DeepTeal,
  pdftitle={论文带读},pdfauthor={论文带读}}
\captionsetup{font=small,labelfont={bf,color=DeepTeal},textfont={color=Ink},skip=5pt}
\setlength{\parindent}{2em}
\setlength{\parskip}{0.28em}
\linespread{1.12}

\titleformat{\section}{\Large\bfseries\color{DeepTeal}}{}{0pt}{}
\titlespacing*{\section}{0pt}{0pt}{6pt}
\titleformat{\subsection}{\large\bfseries\color{DeepTeal}}{}{0pt}{}
\titlespacing*{\subsection}{0pt}{5pt}{3pt}

\pagestyle{fancy}
\fancyhf{}
\lhead{\small\color{DeepTeal}\textbf{文献阅读汇报}\quad <论文主题>}
\rhead{\small\color{gray}<作者>, <年>}
\cfoot{\small\color{gray}\thepage\,/\,\pageref{LastPage}}
\renewcommand{\headrulewidth}{0.5pt}
\renewcommand{\headrule}{\hbox to\headwidth{\color{DeepTeal}\leaders\hrule height \headrulewidth\hfill}}

% ---------- 版式组件 ----------
\newcommand{\key}[1]{\textcolor{DeepTeal}{\textbf{#1}}}
\newcommand{\highlight}[1]{\textcolor{WarmOrange}{\textbf{#1}}}
\newcommand{\source}[1]{\par\vspace{1pt}\raggedright\scriptsize\color{gray}\textit{来源：#1}\par}
\newcommand{\sectiontag}[1]{\noindent\colorbox{PaleTeal}{\parbox{\dimexpr\linewidth-2\fboxsep\relax}{\small\color{DeepTeal}\textbf{#1}}}\par\vspace{5pt}}

% 本节路线（浅青绿框）
\newcommand{\roadmap}[1]{\begin{tcolorbox}[colback=PaleTeal,colframe=MidTeal,boxrule=.5pt,arc=1.5mm,left=2.2mm,right=2.2mm,top=1.4mm,bottom=1.4mm]\small\textbf{本节路线：}\quad #1\end{tcolorbox}}
% 大白话（浅蓝框）
\newcommand{\plain}[1]{\begin{tcolorbox}[colback=PaleBlue,colframe=DeepTeal,boxrule=.5pt,arc=1.5mm,left=2.2mm,right=2.2mm,top=1.5mm,bottom=1.5mm]\small\textbf{大白话}\quad #1\end{tcolorbox}}
% 专业定义 + 本文中的位置（白底青绿框）
\newcommand{\definition}[2]{\begin{tcolorbox}[colback=white,colframe=MidTeal,boxrule=.45pt,arc=1.5mm,left=2.2mm,right=2.2mm,top=1.4mm,bottom=1.4mm]\small\textbf{专业定义：}#1\par\vspace{1mm}\textbf{本文中的位置：}#2\end{tcolorbox}}
% 不要这样理解（暖橙框）
\newcommand{\warning}[1]{\begin{tcolorbox}[colback=WarmOrange!10,colframe=WarmOrange,boxrule=.5pt,arc=1.5mm,left=2.2mm,right=2.2mm,top=1.4mm,bottom=1.4mm]\small\textbf{不要这样理解：}\quad #1\end{tcolorbox}}
% 阅读要点（青绿框）
\newcommand{\takeaway}[1]{\begin{tcolorbox}[colback=PaleTeal,colframe=MidTeal,boxrule=.55pt,arc=1.5mm,left=2.2mm,right=2.2mm,top=1.5mm,bottom=1.5mm]\small\textbf{阅读要点}\quad #1\end{tcolorbox}}
% 注意（白底暖橙框）
\newcommand{\noteBox}[1]{\begin{tcolorbox}[colback=white,colframe=WarmOrange,boxrule=.55pt,arc=1.5mm,left=2.2mm,right=2.2mm,top=1.3mm,bottom=1.3mm]\small #1\end{tcolorbox}}
% 论文图（紧跟图注）
\newcommand{\paperfig}[3]{\begin{figure}[H]\centering\includegraphics[width=#2\linewidth]{#1}\caption{#3}\end{figure}}
% 代码文件标注
\newcommand{\codefile}[2]{\noindent\scriptsize\color{gray}\texttt{#1}\quad\textit{#2}\par\vspace{2pt}}

% ---------- 代码样式 ----------
\lstdefinestyle{pythonstyle}{language=Python,basicstyle=\ttfamily\scriptsize\color{Ink},
  keywordstyle=\color{DeepTeal}\bfseries,commentstyle=\color{gray}\itshape,
  stringstyle=\color{WarmOrange},backgroundcolor=\color{CodeBg},frame=single,
  rulecolor=\color{CodeRule},framerule=.35pt,breaklines=true,showstringspaces=false,
  columns=fullflexible,aboveskip=3pt,belowskip=3pt,xleftmargin=1pt,xrightmargin=1pt}

% ---------- 表格列类型 ----------
\newcolumntype{Y}{>{\RaggedRight\arraybackslash}X}
\newcolumntype{C}{>{\Centering\arraybackslash}X}

% ---------- 安全表格（避免 p{}+Y 混用溢出；见 SKILL.md 陷阱 6） ----------
% 用法：
%   \begin{center}\small
%   \begin{tabular}{@{}p{0.26\linewidth}@{\hspace{4mm}}p{0.62\linewidth}@{}}
%   ...\end{tabular}
%   \end{center}
